%% ============================================================
%%  When Does Orthogonality Fail to Protect?
%%  Ridge versus Lasso as Nuisance Estimators in
%%  Double Machine Learning
%%
%%  Group 5 — FEB23020 — Erasmus School of Economics
%%  Supervisor: Anastasija Tetereva
%%
%%  Compile: pdflatex paper_final.tex  (twice for cross-refs)
%%           bibtex paper_final
%%           pdflatex paper_final.tex  (twice more)
%%  Requires: figuresfinal/ folder, eur.png, MyReferencesFile.bib
%% ============================================================
\documentclass[a4paper,10.5pt]{article}

\author{Daniel Garcia Bru\~{n}a, Renske Gosselaar,
        Johnny Hu, Emile Moyson}
\title{When Does Orthogonality Fail to Protect?\\
       Ridge versus Lasso as Nuisance Estimators\\
       in Double Machine Learning}

\newcommand{\program}{Seminar in Machine Learning}
\newcommand{\supervisor}{Anastasija Tetereva}

%% --- Packages (deduplicated) --------------------------------
\usepackage[british]{babel}
\usepackage[onehalfspacing]{setspace}
\usepackage[margin=2.5cm]{geometry}
\usepackage[T1]{fontenc}
\usepackage[utf8]{inputenc}
%%\usepackage{lmodern}        %% uncomment if available on your system
%%\usepackage{microtype}      %% uncomment if available on your system
\usepackage{graphicx,booktabs}
\usepackage{natbib}
\usepackage{amsmath,amsfonts,amssymb,amsthm}
\usepackage{hyperref}
\usepackage{enumitem}
\usepackage{array}
\usepackage{xcolor}
\usepackage{colortbl}
\usepackage{float}
\usepackage{caption}
\usepackage{subcaption}
\usepackage{multirow}
\usepackage{algpseudocode}
\usepackage{algorithm}

%% --- Theorem environments -----------------------------------
\newtheorem{conjecture}{Conjecture}
\newtheorem{remark}{Remark}
\newtheorem{definition}{Definition}

%% --- Colour definitions -------------------------------------
\definecolor{headerblue}{HTML}{1A5F9E}
\definecolor{lightblue}{HTML}{E6F1FB}
\definecolor{lightgreen}{HTML}{EAF3DE}
\definecolor{lightred}{HTML}{FCEBEB}
\definecolor{mutedgray}{HTML}{5F5E5A}

%% --- Custom table column types ------------------------------
\newcolumntype{C}[1]{>{\centering\arraybackslash}p{#1}}
\newcolumntype{L}[1]{>{\raggedright\arraybackslash}p{#1}}

%% --- Convenience macros -------------------------------------
\newcommand{\E}{\mathbb{E}}
\newcommand{\norm}[1]{\left\lVert #1 \right\rVert}
\newcommand{\Ltwo}{L_2(P)}
\newcommand{\thetahat}{\hat{\theta}}
\newcommand{\thetazero}{\theta_0}
\newcommand{\hatg}{\hat{g}}
\newcommand{\hatm}{\hat{m}}

%% --- Caption style ------------------------------------------
\captionsetup{font=small, labelfont=bf, margin=10pt,
              skip=6pt, justification=justified}

\pagestyle{plain}

\begin{document}

%% ---- Title page --------------------------------------------
\begin{titlepage}
\makeatletter
\begin{center}
    \textsc{Erasmus University Rotterdam}
    \par \textsc{Erasmus School of Economics}

    \vfill \hrule height .08em \bigskip
    \par\huge\@title\bigskip
    \par\Large\@author\\
    \hrule height .08em\normalsize

    \vfill
    
\includegraphics[width=\textwidth,height=0.15\textheight,
                     keepaspectratio]{eur}
    \vfill

    \begin{tabular}{ll}
        \toprule
        Supervisor: & \supervisor \\
        Programme:  & \program    \\
        Date final version: & \@date \\
        \bottomrule
    \end{tabular}

    \vfill
\end{center}
\makeatother
\end{titlepage}

\newpage

%% ---- Abstract ----------------------------------------------
\begin{abstract}
\noindent
Double Machine Learning (DML) protects causal estimates from first-stage estimation errors by means of Neyman's orthogonality condition. In practice, despite the availability of many estimator choices, Lasso is the near-universal default. This paper argues that orthogonality should be understood as a finite insurance mechanism rather than an unconditional guarantee, and that this default choice is not innocuous. Motivated by \citet{shenxiu2025}, who show that Lasso asymptotically collapses to a zero predictor under dense weak signals, the regime that recent evidence suggests characterises most economic data, we investigate whether this prediction failure propagates into invalid causal inference. In a 648-cell Monte Carlo simulation, we find that it does: Lasso-based DML coverage drops to 35\% at empirically realistic signal strengths, while Ridge-based DML maintains 69\%. The mechanism is direct. Lasso's thresholding discards genuine but individually weak coefficients, inflating the nuisance error product beyond the orthogonality budget. Counterintuitively, larger samples worsen the problem rather than alleviating it. Semi-synthetic validation on FRED-MD and Barro-Lee predictor data confirms these findings, and asymmetric designs reveal that the breakdown requires dense weak signals in both nuisance functions simultaneously. When either equation has a dominant sparse component, Lasso is protected. Since Ridge fails under sparse signals, neither method dominates. We propose a simple diagnostic: comparing out-of-fold $R^2$ from both methods on each nuisance regression identifies the researcher's exposure and guides the choice.
\end{abstract}
\newpage


%% ============================================================
\section{Introduction}\label{sec:intro}
%% ============================================================

Researchers in economics increasingly face high-dimensional settings where many potential confounders must be controlled for simultaneously. Double Machine Learning (DML) \citep{chernozhukov2018}, provides an attractive framework for this problem. It uses flexible machine learning methods to estimate nuisance functions while preserving valid inference on the causal parameter of interest. The key theoretical guarantee is Neyman orthogonality: first-stage estimation errors affect the causal estimate only at second order, meaning that even imperfect nuisance estimates can support valid confidence intervals provided they converge at a sufficiently fast rate. This property has made DML one of the most widely adopted tools for causal inference in high-dimensional econometrics.

In the vast majority of applied DML implementations, the nuisance learner is Lasso. This choice is so standard that it is rarely justified explicitly. Lasso is the default in the \texttt{DoubleML} packages for both R and Python, in textbook treatments, and in most published applications. The implicit assumption underlying this choice is that the true nuisance functions depend on a sparse set of predictors: a few variables matter, most do not. Under this assumption, Lasso's ability to perform variable selection makes it the natural candidate. However, this assumption is embedded in the estimation procedure rather than tested against the data, and its validity determines whether the resulting inference can be trusted.

Recent theoretical work challenges this assumption directly. \citet{shenxiu2025} demonstrate that under dense weak signals, where many predictors each carry a small amount of information, Lasso performs no better than predicting zero. Its thresholding mechanism, designed to eliminate irrelevant variables, instead discards true but individually weak coefficients. On the other hand, Ridge regression, which shrinks all coefficients proportionally without eliminating any, continues to extract useful signal in this regime. Crucially, \citet{giannone2021} and \citet{gu2020} provide evidence that economic and financial data typically exhibits precisely this dense weak signal structure: predictability is distributed across many indicators, not concentrated in a few. If this characterisation is correct, then the default nuisance learner in DML is fundamentally mismatched to the data it is applied to.

This raises a natural question for causal inference: does Lasso's prediction failure under dense weak signals compromise DML inference? The answer is not obvious. Orthogonality is specifically designed to absorb first-stage errors. The relevant condition is that the product of the nuisance estimation errors, $\sqrt{n}\cdot\Delta_n$, must vanish. If orthogonality provides a sufficiently large buffer, Lasso's prediction failure might be absorbed without consequence for inference. If the buffer is finite and Lasso's errors exceed it, confidence intervals become invalid. This paper tests which of these two scenarios holds.

We find that orthogonality is a finite buffer that is exhausted under dense weak signals. In a 648-cell Monte Carlo simulation systematically varying signal strength, sparsity, dimensionality, predictor correlation, and sample size, DML-Lasso coverage dropped to 35\% at signal strengths corresponding to the 25th percentile of published economics papers, while DML-Ridge maintained 69\%. The mechanism is direct: Lasso's inability to detect dense weak coefficients produces poor nuisance estimates, which inflate $\sqrt{n}\cdot\Delta_n$ beyond the orthogonality budget, producing confidence intervals that exclude the true parameter 65\% of the time. The framing of orthogonality as a finite insurance mechanism, rather than an unconditional guarantee, is central to understanding this result. Like any insurance policy, orthogonality has a coverage limit. When the claim (the nuisance estimation error) exceeds that limit, the protection fails.

The recommendation, however, is not to replace Lasso with Ridge unconditionally. Under sparse strong signals, where a few predictors dominate, Ridge's uniform shrinkage produces irreducible bias that generates its own coverage breakdown. The largest gap in the entire simulation grid was 91 percentage points favouring Lasso (Lasso 94\%, Ridge 3\%) in the sparse high-dimensional regime. Neither method dominates across all signal environments. The question is not which method is better in general, but which method is better for the data at hand.

Semi-synthetic designs on FRED-MD predictor data revealed that the coverage breakdown requires both the treatment and the outcome equation to involve dense weak signals simultaneously. When either equation contains a dominant sparse component, Lasso exploits the compositional structure of the outcome nuisance function and coverage is preserved.

Counterintuitively, increasing sample size widened the coverage gap from 9 percentage points at $n = 300$ to 27 at $n = 500$ to 34 at $n = 1000$. This occurs because Lasso's prediction error does not improve with $n$ under dense signals \citep{shenxiu2025}, while the orthogonality budget $n^{-1/2}$ shrinks. A researcher who assumes that a large sample protects against first-stage estimation errors is wrong in precisely this setting.

The practical contribution of this paper is a simple diagnostic procedure. Before committing to a nuisance learner, the researcher should run both Lasso and Ridge on each nuisance regression and compare the out-of-fold $R^2$. If Ridge substantially outperforms Lasso on both regressions, the data environment is likely characterised by dense weak signals and Ridge should be preferred. If the gap appears on only one regression, the asymmetric structure provides partial protection. If both methods achieve similar $R^2$, the choice does not materially affect inference. This diagnostic requires no knowledge of the true nuisance functions and avoids the ambiguities associated with the product-based $\sqrt{n}\cdot\Delta_n$ measure, which we show is an upper bound via the Cauchy-Schwarz inequality that can be misleading when the two nuisance functions have different sparsity structures.

The findings are validated on two real predictor datasets. FRED-MD \citep{mccracken2016}, consisting of 122 US macroeconomic indicators, serves as the predictor source for the semi-synthetic exercise, where Ridge produced a strictly lower nuisance error product than Lasso in all 12 symmetric dense cells with ratios stable at 1.41 to 1.60. The Barro-Lee cross-country growth dataset \citep{barro1994}, with 61 predictors and $p/n = 0.678$, extends the validation to a higher dimensionality ratio and a markedly different correlation structure (mean $|\rho_{jk}| = 0.452$), where Ridge maintained an 8 percentage point coverage advantage despite the high predictor correlation that the pure simulation showed typically compresses the gap.

The remainder of the paper is organised as follows. Section~\ref{sec:literature} reviews the relevant literature. Section~\ref{sec:framework} presents the econometric framework. Section~\ref{sec:design} describes the simulation design. Section~\ref{sec:results} presents the results. Section~\ref{sec:implications} discusses practical implications. Section~\ref{sec:conclusion} concludes.


%% ============================================================
\section{Literature Review}\label{sec:literature}
%% ============================================================

This paper connects three strands of literature that have developed largely in isolation: the DML framework for causal inference, the prediction theory under dense weak signals, and the emerging evidence on the sensitivity of causal inference to first-stage estimation quality. We review each strand and identify the gap that motivates this study.


\subsection{Double Machine Learning and the Nuisance Learner Choice}\label{sec:lit_dml}

The DML framework was introduced by \citet{chernozhukov2018} as a general approach to estimating low-dimensional causal parameters in the presence of high-dimensional nuisance functions. The key innovation is the combination of Neyman orthogonality, which ensures that first-stage estimation errors enter the causal estimate only at second order, with cross-fitting, which eliminates the overfitting bias that arises when the same observations are used to estimate the nuisance functions and the causal parameter. If regularity conditions hold and the nuisance estimates converge at a sufficiently fast rate, then the DML estimator is $\sqrt{n}$-consistent, asymptotically normal, and admits valid confidence intervals.

The theoretical foundations of Lasso-based inference in high-dimensional settings were laid by \citet{belloni2014}, who developed post-double-selection methods for treatment effects and instrumental variables. Their framework relies explicitly on approximate sparsity: the assumption that nuisance functions can be well-approximated by a linear combination of a small number of predictors. Under this assumption, Lasso achieves the convergence rate needed for the product condition $\sqrt{n}\cdot\Delta_n \to 0$, and thus valid inference follows. This line of work was later extended to a range of causal parameters including average treatment effects, local average treatment effects, and heterogeneous treatment effects \citep{belloni2014, semenova2021}.

Despite the theoretical generality of DML, which permits any sufficiently accurate machine learning method as the nuisance learner, in practice Lasso is almost exclusively the choice. The \texttt{DoubleML} packages for R \citep{bach2022} and Python use Lasso as the default, and the majority of published applications follow this convention. Alternative learners such as Ridge regression, random forests, boosting, and neural networks are permitted by the theory but are rarely used in practice. This concentration on a single method is consequential because the theoretical guarantees of DML are conditional on the nuisance learner achieving a sufficient convergence rate, which in turn depends on the match between the learner's implicit assumptions and the true data-generating process.

The question of which nuisance learner to use has received surprisingly little attention. \citet{chernozhukov2018} note that the choice should be guided by the structure of the nuisance functions but do not provide concrete guidance for distinguishing between candidate learners in practice. \citet{knaus2021} compares several machine learning methods within DML for heterogeneous treatment effects but focuses on estimation accuracy rather than inference validity. To our knowledge, no prior work has systematically studied how the choice between Lasso and Ridge as nuisance learners affects the coverage of DML confidence intervals across different signal environments.


\subsection{Prediction Under Dense Weak Signals}\label{sec:lit_prediction}

The distinction between sparse and dense signal structures has deep roots in the statistics literature. Lasso \citep{tibshirani1996} was designed for the sparse regime, where variable selection is both feasible and beneficial. Its theoretical properties, including consistency in variable selection and optimal prediction rates, are established under conditions that require approximate sparsity \citep{buhlmann2011, bickel2009}. Ridge regression \citep{hoerl1970} takes a fundamentally different approach. Instead it retains all predictors and then applies uniform shrinkage proportional to each coefficient's magnitude. This makes Ridge suboptimal when the true model is sparse, as it uses regularisation capacity on irrelevant predictors, but advantageous when many predictors each carry a small amount of signal.

The theoretical comparison between the two approaches was sharpened by \citet{shenxiu2025}, who derived minimax prediction risk bounds for both estimators as a function of signal strength and sparsity. Their main result is that under dense weak signals, Lasso's prediction risk does not converge to zero, even asymptotically. In the limit, Lasso collapses to predicting zero because its thresholding mechanism cannot distinguish genuine weak coefficients from noise. On the other hand, Ridge achieves a prediction risk that decreases with sample size, as its proportional shrinkage preserves the aggregate signal carried by many small coefficients. The boundary between the two regimes depends on the ratio of sparsity to dimensionality and on the signal-to-noise ratio.

This theoretical distinction would be of limited practical relevance if economic data were predominantly sparse. However, evidence suggests otherwise. \citet{giannone2021} study the predictability of macroeconomic variables using a large panel of indicators and conclude that prediction is driven by many small contributions rather than a few large ones. \citet{gu2020} reach a similar conclusion in the context of asset pricing, finding that the cross-section of expected returns is better captured by methods that accommodate dense signals. In both cases, Ridge and related shrinkage estimators outperform Lasso in terms of out-of-sample prediction accuracy.

These findings have direct implications for any downstream task that relies on prediction as an intermediate step. DML is precisely such a task: the quality of the causal estimate depends on the quality of the nuisance predictions. If the nuisance functions exhibit the dense weak signal structure documented by \citet{giannone2021} and \citet{gu2020}, then the default choice of Lasso may be actively harmful, not merely suboptimal. This paper tests whether this theoretical concern translates into a measurable consequence for inference.


\subsection{Sensitivity of Causal Inference to First-Stage Quality}\label{sec:lit_sensitivity}

The theoretical conditions under which DML delivers valid inference are well understood. The product rate condition $\|\hat{\ell} - \ell_0\|_{L_2} \cdot \|\hat{m} - m_0\|_{L_2} = o_P(n^{-1/2})$ requires that the product of the two nuisance estimation errors converges faster than $n^{-1/2}$. This condition is weaker than requiring either nuisance estimate to converge at the $n^{-1/4}$ rate individually, and it is this product structure that gives DML its flexibility: one well-estimated nuisance function can partially compensate for a poorly estimated one.

However, the practical robustness of this condition has been questioned. \citet{chernozhukov2022} discuss the sensitivity of DML estimates to the quality of nuisance learners and emphasise that the rate condition is a necessary requirement, not merely a sufficient one. When it fails, the DML estimator is biased and confidence intervals are too narrow. The failure is not detectable from the output of the procedure itself, which continues to produce point estimates and standard errors regardless of whether the underlying rate condition holds.

A related concern is raised by the literature on the fragility of sparsity assumptions in econometric practice. \citet{giannone2021} note that the sparsity assumption underlying Lasso-based methods is difficult to verify in practice and may be violated in precisely the settings where high-dimensional methods are most needed. This echoes the broader point made by \citet{leeb2005}, who show that inference after model selection can be severely distorted when the selected model omits relevant variables, even if the omitted variables have individually small coefficients.

The connection between prediction quality and inference validity in DML has not, to our knowledge, been studied through the lens of dense versus sparse signal structures. Existing simulation studies of DML \citep{bach2022, chernozhukov2018} typically use sparse data-generating processes that are favourable to Lasso, and do not explore the behaviour of the estimator when the sparsity assumption is violated. Our contribution fills this gap by systematically varying the signal environment from sparse strong to dense weak and documenting the consequences for inference validity under both Lasso and Ridge as nuisance learners.


%% ============================================================
\section{Econometric Framework}\label{sec:framework}
%% ============================================================

This section formalises the intuition developed in the introduction. We present the partially linear model and the DML estimator (Section~\ref{sec:plr}), reinterpret the rate condition as a finite orthogonality budget (Section~\ref{sec:buffer}), explain why the choice between Lasso and Ridge determines whether this budget is exhausted (Section~\ref{sec:signal_structure}), and state the conjecture that motivates the simulation design (Section~\ref{sec:conjecture}).


\subsection{The Partially Linear Model and DML Estimator}\label{sec:plr}

Consider the partially linear regression (PLR) model
\begin{align}
    Y_i &= \theta_0 D_i + g_0(X_i) + \varepsilon_i, \quad \mathbb{E}[\varepsilon_i \mid X_i, D_i] = 0, \label{eq:plr_y} \\
    D_i &= m_0(X_i) + v_i, \quad \mathbb{E}[v_i \mid X_i] = 0, \label{eq:plr_d}
\end{align}
where $Y_i$ is the outcome, $D_i$ is the treatment, $X_i \in \mathbb{R}^p$ is a vector of controls, $\theta_0$ is the causal parameter of interest, $g_0(X) = \mathbb{E}[Y \mid X] - \theta_0 \mathbb{E}[D \mid X]$ is the direct effect of the controls on the outcome, and $m_0(X) = \mathbb{E}[D \mid X]$ is the conditional mean of the treatment given the controls. The functions $g_0$ and $m_0$ are nuisance parameters: they must be estimated to recover $\theta_0$, but they are not of primary interest.

Define the outcome nuisance function as $\ell_0(X) = \mathbb{E}[Y \mid X] = \theta_0 m_0(X) + g_0(X)$. Note that $\ell_0$ inherits structure from both $m_0$ and $g_0$ through the identity $\ell_0 = \theta_0 m_0 + g_0$. This compositional structure will play a central role in the asymmetric analysis of Section~\ref{sec:asymmetric}.

The DML estimator of $\theta_0$ proceeds in three steps. First, the sample is split into $K$ folds. For each fold $k$, the nuisance functions $\hat{\ell}$ and $\hat{m}$ are estimated on the complement of fold $k$ using a machine learning method. Second, the out-of-fold residuals are formed:
\begin{equation}
    \hat{W}_i = Y_i - \hat{\ell}(X_i), \quad \hat{V}_i = D_i - \hat{m}(X_i).
\end{equation}
Third, the causal parameter is estimated by the cross-fitted moment condition:
\begin{equation}\label{eq:dml_estimator}
    \hat{\theta} = \frac{\sum_{i=1}^{n} \hat{V}_i \hat{W}_i}{\sum_{i=1}^{n} \hat{V}_i^2}.
\end{equation}
The variance is estimated by the sandwich formula:
\begin{equation}\label{eq:variance}
    \hat{\sigma}^2_\theta = \frac{n^{-1} \sum_{i=1}^{n} \hat{V}_i^2 \hat{\varepsilon}_i^2}{\left( n^{-1} \sum_{i=1}^{n} \hat{V}_i^2 \right)^2},
\end{equation}
where $\hat{\varepsilon}_i = \hat{W}_i - \hat{\theta} \hat{V}_i$. The 95\% confidence interval is $\hat{\theta} \pm 1.96 \cdot \hat{\sigma}_\theta / \sqrt{n}$.

To reduce sensitivity to a single random fold partition, the cross-fitting procedure is repeated $S$ times with independent random splits. The final estimate is the average across repetitions:
\begin{equation}
    \tilde{\theta} = \frac{1}{S} \sum_{s=1}^{S} \hat{\theta}^{(s)},
\end{equation}
with variance
\begin{equation}
    \tilde{\sigma}^2_\theta = \frac{1}{S} \sum_{s=1}^{S} \left[ \hat{\sigma}^{(s)2}_\theta + \left(\hat{\theta}^{(s)} - \tilde{\theta}\right)^2 \right],
\end{equation}
which accounts for both the within-split estimation uncertainty and the across-split variation due to different fold assignments.


\subsection{Orthogonality as a Finite Buffer}\label{sec:buffer}

The key property of the DML estimator is Neyman orthogonality. A Taylor expansion of the moment condition around the true nuisance functions yields
\begin{equation}\label{eq:expansion}
    \sqrt{n}(\hat{\theta} - \theta_0) = \underbrace{\frac{n^{-1/2} \sum_i V_i \varepsilon_i}{n^{-1} \sum_i V_i^2}}_{\text{influence function}} + \underbrace{\sqrt{n} \cdot R_n}_{\text{remainder}},
\end{equation}
where the influence function term is asymptotically normal by the central limit theorem, and $R_n$ is a remainder that depends on the nuisance estimation errors. Orthogonality ensures that the first-order terms in $R_n$ vanish, leaving only second-order contributions.

The remainder is bounded by
\begin{equation}\label{eq:remainder_bound}
    |R_n| \;\leq\; \|\hat{\ell} - \ell_0\|_{L_2} \cdot \|\hat{m} - m_0\|_{L_2} \;=:\; \Delta_n,
\end{equation}
via the Cauchy-Schwarz inequality, where $\|f\|_{L_2} = \left(n^{-1} \sum_{i=1}^{n} f(X_i)^2\right)^{1/2}$ denotes the empirical $L_2$ norm. The rate condition for valid inference is
\begin{equation}\label{eq:rate_condition}
    \sqrt{n} \cdot \Delta_n = \sqrt{n} \cdot \|\hat{\ell} - \ell_0\|_{L_2} \cdot \|\hat{m} - m_0\|_{L_2} \;\to_P\; 0.
\end{equation}
When this condition holds, the remainder is asymptotically negligible and $\hat{\theta}$ is $\sqrt{n}$-consistent with valid confidence intervals.

We find it useful to interpret this condition as a \emph{finite budget}. The orthogonality mechanism provides protection of magnitude $n^{-1/2}$: it absorbs nuisance errors up to this threshold. The nuisance estimation errors generate a ``claim'' of size $\Delta_n$ against this budget. If the claim is smaller than the budget ($\sqrt{n} \cdot \Delta_n \to 0$), the insurance holds and inference is valid. If the claim exceeds the budget ($\sqrt{n} \cdot \Delta_n \not\to 0$), the protection is exhausted and the confidence interval is too narrow.

This framing highlights two distinct sources of failure. The budget $n^{-1/2}$ shrinks as the sample grows, which should help inference by reducing the influence of individual observations. But the claim $\Delta_n$ depends on the quality of the nuisance estimates, which depends on the signal structure and the estimation method. If the estimation method produces errors that do not shrink with $n$, the budget contracts while the claim remains constant, and $\sqrt{n} \cdot \Delta_n$ diverges. This is precisely what happens when Lasso is applied to dense weak signals.

It is important to note that $\Delta_n$ is a Cauchy-Schwarz \emph{upper bound} on the actual remainder $|R_n|$, not the remainder itself. The actual remainder depends on the inner product $n^{-1}\sum_i (\hat{\ell}_i - \ell_{0,i})(\hat{m}_i - m_{0,i})$, which can be much smaller than the product of norms when the two error vectors are misaligned. We return to this distinction in Section~\ref{sec:asymmetric}, where it explains why $\sqrt{n}\cdot\Delta_n$ tracks coverage tightly in symmetric designs but can be misleading under asymmetric nuisance structures.


\subsection{Why Signal Structure Determines Nuisance Learner Failure}\label{sec:signal_structure}

In our simulation design, the nuisance functions are linear: $m_0(X) = X\beta_m$ and $g_0(X) = X\beta_g$. The signal structure is governed by two parameters for each equation: the signal strength $\tau$, defined as $\|\beta\|^2 = \sum_{j=1}^p \beta_j^2$, and the sparsity $s/p$, the fraction of nonzero coefficients. Together, these determine the population $R^2$ of the nuisance regression:
\begin{equation}
    R^2 = \frac{\tau}{\tau + 1},
\end{equation}
in the standardised design where $X$ has identity covariance and the noise has unit variance. The coefficient vector is constructed as
\begin{equation}\label{eq:beta_construction}
    \beta_j = \begin{cases} \sqrt{\tau/s} & \text{if } j \leq s, \\ 0 & \text{if } j > s, \end{cases}
\end{equation}
where $s = \lfloor (s/p) \cdot p \rfloor$. Each active coefficient has magnitude $\sqrt{\tau/s}$, so that increasing density (higher $s/p$) at fixed $\tau$ distributes the same total signal across more coefficients, making each individual coefficient smaller.

Lasso and Ridge respond to this structure in opposite ways. Lasso \citep{tibshirani1996} solves
\begin{equation}
    \hat{\beta}^{\text{Lasso}} = \argmin_\beta \|Y - X\beta\|^2_2 + \lambda \|\beta\|_1,
\end{equation}
where the $\ell_1$ penalty induces exact zeros: coefficients below a data-dependent threshold are set to zero. When the true model is sparse ($s/p$ small) and individual coefficients are large ($\sqrt{\tau/s}$ large), this thresholding correctly eliminates the irrelevant predictors and retains the important ones. But when $s/p$ is large and individual coefficients are small, the genuine coefficients fall below the threshold and are discarded. As \citet{shenxiu2025} show, in the dense weak limit, Lasso's prediction collapses to the zero vector: it eliminates all signal because no individual coefficient is detectable.

Ridge \citep{hoerl1970} solves
\begin{equation}
    \hat{\beta}^{\text{Ridge}} = \argmin_\beta \|Y - X\beta\|^2_2 + \lambda \|\beta\|^2_2,
\end{equation}
where the $\ell_2$ penalty shrinks all coefficients toward zero proportionally, without eliminating any. This is suboptimal when the true model is sparse, because Ridge wastes regularisation capacity on null predictors. But under dense weak signals, Ridge's proportional shrinkage preserves the aggregate signal: each coefficient is attenuated but not eliminated, and the sum of many small attenuated contributions still produces useful predictions. The prediction risk of Ridge converges to zero at a rate that depends on $p/n$ and $\tau$ but does not require individual coefficients to exceed a detection threshold.

The consequence for DML is immediate. Under dense weak signals, Lasso produces nuisance estimates with elevated $L_2$ error, leading to a large $\Delta_n$. Under sparse strong signals, Ridge produces nuisance estimates with elevated $L_2$ error, for the opposite reason. The rate condition~\eqref{eq:rate_condition} is method-dependent: the same data can satisfy the condition under one method and violate it under another.


\subsection{The Conjecture}\label{sec:conjecture}

The preceding analysis motivates the following conjecture.

\medskip
\noindent\textbf{Conjecture.} \textit{As signal strength $\tau$ decreases and density $s/p$ increases, Lasso's nuisance error product $\Delta_n^{\text{Lasso}}$ will exceed Ridge's $\Delta_n^{\text{Ridge}}$ by a growing margin, eventually violating the rate condition and producing DML confidence intervals with coverage substantially below the nominal 95\% level. The coverage gap between DML-Ridge and DML-Lasso will be largest at intermediate signal strengths ($\tau \in [0.05, 0.2]$), will widen with dimensionality ($p/n$), will narrow with predictor correlation ($\rho$), and will widen with sample size ($n$) because Lasso's prediction error does not improve with $n$ under dense signals while the orthogonality budget $n^{-1/2}$ shrinks.}
\medskip

We extend this conjecture to asymmetric nuisance structures. Since $\Delta_n$ is the product of two nuisance errors, a single poorly estimated function should in principle be sufficient to violate the rate condition. However, the product structure also implies that one well-estimated function may partially compensate for a poorly estimated one. We conjecture that when either $\beta_m$ or $\beta_g$ contains a dominant sparse component, this component flows into $\ell_0 = \theta_0 m_0 + g_0$ and provides Lasso with an exploitable sparse structure, preserving coverage even if the other component is dense and weak.

The simulation design described in the next section is constructed to test this conjecture across a comprehensive grid of signal environments.


%% ============================================================
\section{Research Design}\label{sec:design}
%% ============================================================

The empirical analysis comprises three components, each designed to address a distinct inferential challenge. The pure simulation (Section~\ref{sec:pure_sim}) uses fully synthetic data with known nuisance functions to evaluate bias, coverage, and the nuisance error product across a comprehensive grid of signal environments. The semi-synthetic exercise (Section~\ref{sec:semi_synth}) embeds synthetic treatments and outcomes in real predictor matrices to test whether the mechanism documented in the pure simulation survives under empirical covariance structures. The empirical diagnostic exercise (Section~\ref{sec:empirical_design}) applies the proposed diagnostic to datasets where the true treatment effect is unknown, demonstrating its feasibility in practice. Section~\ref{sec:metrics} defines the evaluation metrics that are common to all three components.


\subsection{Pure Simulation}\label{sec:pure_sim}

\subsubsection*{Data-Generating Process}

The data-generating process follows the partially linear model of equations~\eqref{eq:plr_y}--\eqref{eq:plr_d}. For each replication $b = 1, \ldots, B$, the predictor matrix $X^{(b)} \in \mathbb{R}^{n \times p}$ is drawn independently from a multivariate normal distribution $X^{(b)}_i \sim \mathcal{N}(0, \Sigma)$, where $\Sigma$ follows the Toeplitz structure $\Sigma_{jk} = \rho^{|j-k|}$ for $j, k = 1, \ldots, p$. Each column of $X^{(b)}$ is re-standardised to have sample mean zero and sample variance one within the replication. The noise terms $v^{(b)}_i \sim \mathcal{N}(0,1)$ and $\varepsilon^{(b)}_i \sim \mathcal{N}(0,1)$ are drawn independently of each other and of $X^{(b)}$.

The coefficient vectors $\beta_m$ and $\beta_g$ are set equal ($\beta_m = \beta_g = \beta$) in the pure simulation. The common coefficient vector is constructed as in equation~\eqref{eq:beta_construction}, ensuring that the total signal energy $\tau = \|\beta\|^2$ is distributed equally across $s = \lfloor (s/p) \cdot p \rfloor$ active predictors. Increasing the density parameter $s/p$ at fixed $\tau$ reduces the magnitude of each individual coefficient from $\sqrt{\tau/s}$ toward zero. The true treatment effect is set to $\theta_0 = 1$ throughout.

The nuisance functions are $m_0(X) = X\beta$ and $\ell_0(X) = (\theta_0 + 1) \cdot X\beta = 2X\beta$ (since $\beta_m = \beta_g = \beta$ and $\theta_0 = 1$). The population $R^2$ of each nuisance regression, under the standardised design with $\text{Var}(\varepsilon) = 1$, is $R^2 = \tau / (\tau + 1)$.


\subsubsection*{Parameter Space}

The simulation grid spans five dimensions, chosen to capture the key factors identified in the theoretical framework. Table~\ref{tab:calibration} maps signal strength to population $R^2$ and empirical context.

\begin{table}[H]
\centering
\small
\caption{Signal strength calibration: mapping $\tau$ to population $R^2$ and empirical context.}
\label{tab:calibration}
\begin{tabular}{lccl}
\toprule
$\tau$ & $R^2$ (\%) & Empirical context \\
\midrule
5.0  & 83.3 & Strong: structural macro models \\
2.0  & 66.7 & Strong: well-specified treatment equations \\
1.0  & 50.0 & Moderate: typical micro regressions \\
0.5  & 33.3 & Moderate: macro forecasting \\
0.2  & 16.7 & Weak: typical macro panel \\
0.1  & 9.1  & Weak: 25th percentile of published work \\
0.05 & 4.8  & Very weak: equity premium prediction \\
0.01 & 1.0  & Near-zero: noise dominates \\
\bottomrule
\end{tabular}
\end{table}

\textit{Signal strength} ($\tau$) governs the detectability of individual coefficients. The values $\tau \in \{5.0, 2.0, 1.0, 0.5, 0.2, 0.1, 0.05, 0.01\}$ correspond to population $R^2$ values ranging from 83.3\% to 1.0\%. The value $\tau = 0.1$ ($R^2 = 9.1\%$) was calibrated to approximate the 25th percentile of $R^2$ values reported in published economics papers, placing it squarely in the empirically relevant range.

\textit{Sparsity} ($s/p$) determines how the signal is distributed across predictors. The values $s/p \in \{0.05, 0.30, 1.0\}$ represent three qualitatively distinct regimes: highly sparse (5\% of predictors active, corresponding to the assumption embedded in Lasso), moderately sparse (30\%), and fully dense (100\%, the regime studied by \citealt{shenxiu2025}). At fixed $\tau$, moving from $s/p = 0.05$ to $s/p = 1.0$ reduces the individual coefficient magnitude by a factor of $\sqrt{20} \approx 4.5$, making each coefficient substantially harder to detect.

\textit{Dimensionality} ($p$) controls the number of predictors relative to the sample size. The values $p \in \{100, 200, 500\}$ at fixed $n$ yield dimensionality ratios $p/n$ ranging from 0.10 to 1.67.

\textit{Sample size} ($n$) determines both the precision of the nuisance estimates and the magnitude of the orthogonality budget $n^{-1/2}$. The values $n \in \{300, 500, 1000\}$ were chosen to span from moderately sized observational studies to large panel datasets.

\textit{Predictor correlation} ($\rho$) controls the dependence structure. The values $\rho \in \{0.0, 0.5, 0.8\}$ range from independent predictors to high collinearity.

The full factorial design comprised $3 \times 3 \times 8 \times 3 \times 3 = 648$ cells, of which 641 produced numerically stable results. The seven failures occurred at extreme configurations ($p > n$ with strong signals and low correlation) where both methods produced degenerate variance estimates.

\subsubsection*{DML Implementation}

Each cell was replicated $B = 100$ times with $K = 5$-fold cross-fitting repeated $S = 3$ times. DML-Lasso used LassoCV with $\lambda$ selected over 15 log-spaced values from $10^{-4}$ to $10^1$ via $K_{\text{inner}} = 3$-fold cross-validation. DML-Ridge used RidgeCV with $\lambda$ selected over 60 log-spaced values from $10^{-4}$ to $10^6$. Within-fold re-standardisation prevented information leakage. Full implementation details, including the algorithm pseudocode and the rationale for cross-validated rather than theoretical $\lambda$ selection, are provided in Appendix~\ref{app:algorithm}.


\subsection{Semi-Synthetic Design}\label{sec:semi_synth}

The pure simulation provides full control over the data-generating process at the cost of imposing a parametric covariance structure (Toeplitz) on the predictors. The semi-synthetic design relaxes this restriction by taking the predictor matrix $X$ from real economic datasets while generating the treatment and outcome synthetically. This preserves the empirical dependence structure while keeping $\theta_0 = 1$ known by construction. For each replication $b$, only the noise terms are redrawn; $X$ is held fixed. The coefficient vectors $\beta_m$ and $\beta_g$ are now allowed to differ, enabling asymmetric designs. The DML implementation is identical to the pure simulation.

\subsubsection*{FRED-MD}

The FRED-MD dataset \citep{mccracken2016} yielded $n = 757$ observations and $p = 122$ predictors ($p/n = 0.161$) after McCracken-Ng stationarity transformations (Appendix~\ref{app:fredmd}). Subsamples of $n = 400$ ($p/n = 0.305$) and $n = 250$ ($p/n = 0.488$) were drawn using a fixed random seed. The mean absolute pairwise correlation was 0.116, with only 4.4\% of predictor pairs exceeding $|\rho_{jk}| > 0.5$, placing FRED-MD closest to the $\rho = 0.0$ Toeplitz cells.

Four coefficient designs were tested at each sample size:
\begin{enumerate}[label=(\roman*)]
    \item \textit{Symmetric dense} ($\beta_m = \beta_g$, $s/p = 1.0$, $\tau \in \{0.2, 0.1, 0.05, 0.01\}$): 12 cells.
    \item \textit{Asymmetric: sparse $\beta_m$ + dense $\beta_g$} ($\tau_m = 2.0$, $s_m/p = 0.05$; $\tau_g \in \{0.2, 0.1, 0.05, 0.01\}$, $s_g/p = 1.0$): 12 cells.
    \item \textit{Asymmetric: dense $\beta_m$ + sparse $\beta_g$} ($\tau_m \in \{0.1, 0.05\}$, $s_m/p = 1.0$; $\tau_g = 2.0$, $s_g/p = 0.05$): 6 cells.
    \item \textit{Symmetric sparse baseline} ($\beta_m = \beta_g$, $s/p = 0.05$, $\tau = 2.0$): 3 cells.
\end{enumerate}
In total, 33 cells were evaluated, each with $B = 100$ replications.

\subsubsection*{Barro-Lee}

The Barro-Lee cross-country growth dataset \citep{barro1994} comprised $n = 90$ countries with $p = 61$ predictors ($p/n = 0.678$). Its correlation structure was qualitatively different from FRED-MD: the mean absolute pairwise correlation was 0.452, with 44.2\% of pairs exceeding $|\rho_{jk}| > 0.5$. This high correlation, which the pure simulation showed typically compresses the coverage gap, provided a stringent test of whether Ridge's advantage survived under unfavourable conditions. The same four coefficient designs were applied across 11 cells with $B = 50$ replications (coverage SE $\approx 3.1$~pp).


\subsection{Empirical Diagnostics}\label{sec:empirical_design}

The empirical component asks whether a researcher can detect their exposure to the Lasso failure using only observable quantities. We propose a four-step diagnostic: (1)~signal strength positioning via Ridge $R^2$, (2)~separate $R^2$ comparison on each nuisance regression, (3)~approximate $\sqrt{n}\cdot\hat{\Delta}_n$ from out-of-fold residuals, and (4)~estimate stability across cross-fitting splits. The separate comparison of $R^2$ on each regression, rather than a single product-based diagnostic, is motivated by the finding (Section~\ref{sec:mechanism}) that $\sqrt{n}\cdot\Delta_n$ is a Cauchy-Schwarz upper bound that can be loose under asymmetric nuisance structures.

Two datasets were used: Barro-Lee ($n = 90$, $p = 61$) with four treatment variables, and Goyal-Welch \citep{goyal2008} ($n = 68$, $p = 15$) with financial predictors of the equity premium. FRED-MD was excluded to avoid redundancy with the semi-synthetic design.


\subsection{Evaluation Metrics}\label{sec:metrics}

The following metrics are computed for each cell, reported separately for DML-Lasso and DML-Ridge.

\paragraph{Bias and Root Mean Squared Error.} Bias is defined as $\text{Bias} = B^{-1}\sum_{b=1}^{B}(\hat{\theta}_b - \theta_0)$ and RMSE as $\text{RMSE} = [B^{-1}\sum_{b=1}^{B}(\hat{\theta}_b - \theta_0)^2]^{1/2}$. Both are available only in the simulation and semi-synthetic designs where $\theta_0 = 1$ is known.

\paragraph{Coverage.} Empirical coverage is defined as $\text{Coverage} = B^{-1}\sum_{b=1}^{B}\mathbf{1}\{\theta_0 \in \text{CI}_b\}$, where $\text{CI}_b = \hat{\theta}_b \pm 1.96 \cdot \hat{\sigma}_{\theta,b} / \sqrt{n}$ is the 95\% confidence interval from replication $b$. This is the primary metric. The standard error of the coverage estimate is $\sqrt{c(1-c)/B}$, approximately 2.2~pp at $c = 0.95$ and $B = 100$.

\paragraph{Nuisance Error Product.} The quantity $\sqrt{n}\cdot\Delta_n = \sqrt{n} \cdot \|\hat{\ell} - \ell_0\|_{L_2} \cdot \|\hat{m} - m_0\|_{L_2}$ measures the claim against the orthogonality budget. As established in Section~\ref{sec:buffer}, this provides a Cauchy-Schwarz upper bound on the scaled remainder $\sqrt{n}\cdot|R_n|$, tight when errors are aligned and loose when misaligned.

\paragraph{Out-of-Fold $R^2$.} Prediction quality is measured separately for each nuisance regression: $R^2_\ell = 1 - \sum_i(\hat{\ell}_i - \ell_{0,i})^2 / \sum_i(\ell_{0,i} - \bar{\ell}_0)^2$ for the outcome regression, and analogously $R^2_m$ for the treatment regression. The separate reporting of $R^2_\ell$ and $R^2_m$ is motivated by the finding that the product-based $\sqrt{n}\cdot\Delta_n$ can be misleading in asymmetric settings, while separate $R^2$ values provide a more reliable and interpretable diagnostic.



\section{Results}\label{sec:results}
%% ============================================================

This section presents the simulation and semi-synthetic results. Section~\ref{sec:coverage} documents the central finding: a pronounced coverage breakdown in DML-Lasso under dense weak signals. Section~\ref{sec:mechanism} traces the causal chain from prediction failure through the nuisance error product to inference failure, and discusses the relationship between the theoretical diagnostic $\sqrt{n}\cdot\Delta_n$ and the practically computable out-of-fold $R^2$. Section~\ref{sec:sensitivity} examines the moderating role of each simulation parameter. Section~\ref{sec:asymmetric} investigates when the researcher is protected through asymmetric nuisance structures. Section~\ref{sec:validation} validates the findings on real predictor data from FRED-MD and Barro-Lee.


%% ============================================================
\subsection{Coverage Breakdown Under Dense Weak Signals}\label{sec:coverage}
%% ============================================================

Of the 107 cells in the simulation grid satisfying the weak dense condition ($\tau \leq 0.2$, $s/p = 1.0$), Ridge achieved strictly higher coverage than Lasso in 75 cells (70.1\%). Lasso achieved strictly higher coverage in only 15 cells, and the two methods produced identical coverage in 17 cells.

The magnitude of the advantage was strikingly asymmetric. When Ridge outperformed Lasso, the average coverage gap was \textbf{10.3 percentage points} (pp), with a maximum of 41.0~pp. When Lasso outperformed Ridge, the average gap was only \textbf{1.6~pp}, and in 10 of the 15 cells both methods already had coverage above 90\%. Switching from Lasso to Ridge in the weak dense regime thus constitutes a free option: substantial upside when the data environment favours Ridge, negligible downside when it does not.

The most informative single cell was $n = 1000$, $p = 200$, $\tau = 0.1$, $s/p = 1.0$, $\rho = 0.0$, where Lasso coverage was \textbf{35.0\%} and Ridge coverage was \textbf{69.0\%}. At this signal strength ($\tau = 0.1$, corresponding to $R^2 \approx 9\%$), the implied population $R^2$ falls near the 25th percentile of published economics papers, placing this configuration squarely in the empirically relevant range. A researcher using DML-Lasso at this signal strength would report a 95\% confidence interval that contains the true parameter only 35\% of the time.

\begin{figure}[H]
  \centering
  \includegraphics[width=\textwidth]{fig1.pdf}
  \caption{Empirical 95\% CI coverage as a function of signal strength~$\tau$ for DML-Lasso and DML-Ridge at $p = 200$, $s/p = 1.0$, $\rho = 0.0$, across three sample sizes. As $\tau$ decreases from 5.0 toward 0.1, Lasso coverage declines sharply while Ridge declines more slowly, opening a widening gap. At $\tau = 0.01$, both methods recover toward nominal levels because the nuisance functions approach zero. The dashed line marks the nominal 95\% level.}
  \label{fig:coverage_by_n}
\end{figure}

Figure~\ref{fig:coverage_by_n} displays the coverage trajectories across the full signal spectrum for the three sample sizes. Two patterns are immediately apparent. First, as $\tau$ decreased from 5.0 toward 0.1, Lasso coverage declined sharply while Ridge declined more slowly, producing a widening gap that peaked at intermediate signal strengths. Second, at very weak signals ($\tau = 0.01$), both methods recovered toward nominal coverage. This U-shaped pattern reflects the fact that as $\tau \to 0$, the nuisance functions $\ell_0(X) = 2X\beta$ and $m_0(X) = X\beta$ approach zero, reducing the estimation problem to simple OLS on the residuals, which both methods handle adequately.

Table~\ref{tab:gradient} presents the gradient from strong sparse to weak dense signals at a single representative configuration, illustrating the continuous transition from Lasso dominance to Ridge dominance.

\begin{table}[H]
\centering
\small
\caption{Coverage gradient from strong sparse to weak dense signals. Configuration: $n = 1000$, $p = 200$, $\rho = 0.0$. Each row represents one cell ($B = 100$ replications, $S = 3$ cross-fitting repetitions).}
\label{tab:gradient}
\begin{tabular}{llcccccc}
\toprule
& & \multicolumn{2}{c}{Coverage (\%)} & \multicolumn{2}{c}{$\sqrt{n}\cdot\Delta_n$} & \multicolumn{2}{c}{$R^2_\ell$} \\
\cmidrule(lr){3-4} \cmidrule(lr){5-6} \cmidrule(lr){7-8}
$(\tau,\; s/p)$ & Regime & Lasso & Ridge & Lasso & Ridge & Lasso & Ridge \\
\midrule
$(5.0,\; 0.05)$ & Strong sparse    & 94.0 & 89.0 & 2.57 & 11.96 & --- & --- \\
$(2.0,\; 0.05)$ & Strong sparse    & 94.0 & 85.0 & 2.57 & 11.18 & --- & --- \\
$(1.0,\; 0.30)$ & Moderate         & 91.0 & 79.0 & 8.00 & 10.22 & --- & --- \\
$(0.5,\; 1.0)$  & Transition       & 79.0 & 66.0 & 12.85 & 8.64 & 0.698 & 0.785 \\
$(0.2,\; 1.0)$  & Weak dense       & 49.0 & 60.0 & 9.47 & 6.05 & 0.365 & 0.598 \\
$(0.1,\; 1.0)$  & Weak dense       & \textbf{35.0} & \textbf{69.0} & 5.85 & 4.11 & 0.130 & 0.427 \\
$(0.05,\; 1.0)$ & Weak dense       & 66.0 & 85.0 & 3.14 & 2.56 & 0.026 & 0.257 \\
$(0.01,\; 1.0)$ & Near-zero signal & 92.0 & 92.0 & 0.69 & 0.69 & $-$0.058 & $-$0.033 \\
\bottomrule
\end{tabular}
\end{table}

The gradient reveals a clear crossover. At strong sparse signals ($\tau \geq 2.0$, $s/p = 0.05$), Lasso maintained 94\% coverage while Ridge was depressed at 85--89\%. The crossover occurred between $\tau = 0.5$ and $\tau = 0.2$, after which Ridge dominated. At $\tau = 0.01$, both methods converged to 92\% as the nuisance functions effectively vanished. The most severe distortions occurred at intermediate weak signals ($\tau \in [0.05, 0.2]$), not at the weakest signals, consistent with the intuition that coverage is worst when signals are present but difficult to detect.


%% ============================================================
\subsection{The Mechanism: From Prediction Failure to Inference Failure}\label{sec:mechanism}
%% ============================================================

The theoretical prediction was that coverage breakdown would be driven by an elevated nuisance error product $\Delta_n = \|\hat{\ell} - \ell_0\|_{L_2} \cdot \|\hat{m} - m_0\|_{L_2}$. For valid inference, $\sqrt{n}\cdot\Delta_n$ must vanish; when it does not, the orthogonality buffer is exhausted and confidence intervals become too narrow. This section traces the mechanistic chain and discusses its implications for practical diagnostics.

\subsubsection*{The Chain: $R^2 \to \sqrt{n}\cdot\Delta_n \to$ Coverage}

At the cleanest mechanistic cell ($n = 1000$, $p = 200$, $\tau = 0.1$, $s/p = 1.0$, $\rho = 0.0$), Lasso's out-of-fold $R^2_\ell$ was 0.130, indicating that it explained only 13.0\% of the variation in the outcome nuisance function $\ell_0(X)$. Ridge achieved $R^2_\ell = 0.427$, more than three times Lasso's prediction accuracy. This prediction gap fed directly into the nuisance error product: $\sqrt{n}\cdot\Delta_n$ was 5.85 for Lasso versus 4.11 for Ridge. Both values exceeded the theoretical threshold of approximately~1, but Ridge's lower value was associated with substantially better coverage: 69.0\% versus 35.0\%.

\begin{figure}[H]
  \centering
  \includegraphics[width=\textwidth]{fig2.pdf}
  \caption{The mechanistic chain at $n = 1000$, $p = 200$, $s/p = 1.0$, $\rho = 0.0$. Panel~(a): the nuisance error product $\sqrt{n}\cdot\Delta_n$ declined as $\tau$ decreased, with Ridge consistently below Lasso for $\tau \leq 1.0$. Panel~(b): coverage tracked $\sqrt{n}\cdot\Delta_n$ inversely. The dotted line in~(a) marks $\sqrt{n}\cdot\Delta_n = 1$; the dashed line in~(b) marks 95\% nominal coverage.}
  \label{fig:mechanism}
\end{figure}

Figure~\ref{fig:mechanism} displays the chain across the full signal spectrum. In the left panel, both methods' $\sqrt{n}\cdot\Delta_n$ declined as $\tau$ decreased, with Ridge consistently below Lasso for $\tau \leq 1.0$. The two curves converged at $\tau = 0.01$ (both at 0.69), confirming that the difference between methods was driven entirely by the signal estimation stage. In the right panel, coverage tracked the nuisance error product inversely: higher $\sqrt{n}\cdot\Delta_n$ corresponded to lower coverage, with a nonlinear mapping in which small differences in $\sqrt{n}\cdot\Delta_n$ were associated with large differences in coverage at intermediate values.

\begin{figure}[H]
  \centering
  \includegraphics[width=0.65\textwidth]{fig3.pdf}
  \caption{Out-of-fold $R^2_\ell$ for the outcome nuisance regression at $n = 1000$, $p = 200$, $s/p = 1.0$, $\rho = 0.0$. Ridge consistently outperformed Lasso, with the gap widening as signals weakened. At $\tau = 0.01$, both methods produced negative $R^2$ (worse than a constant predictor), confirming that no learnable signal remained.}
  \label{fig:R2}
\end{figure}

Figure~\ref{fig:R2} isolates the first link in the chain. Ridge's $R^2_\ell$ advantage widened as $\tau$ decreased: from a negligible gap at $\tau = 5.0$ (both near 0.97) to a gap of 0.297 at $\tau = 0.1$ (Ridge: 0.427, Lasso: 0.130). At $\tau = 0.01$, both methods produced negative $R^2$ values ($-0.058$ for Lasso, $-0.033$ for Ridge), confirming convergence to the constant predictor at near-zero signal strength.

\begin{table}[H]
\centering
\small
\caption{Mechanistic chain at $n = 1000$, $p = 200$, $s/p = 1.0$. The lower panel shows the effect of correlation at fixed $\tau = 0.1$, demonstrating how $\rho$ moderates the mechanism through improved effective dimensionality.}
\label{tab:mechanism}
\begin{tabular}{lccccccc}
\toprule
& & \multicolumn{2}{c}{$\sqrt{n}\cdot\Delta_n$} & \multicolumn{2}{c}{Coverage (\%)} & \multicolumn{2}{c}{$R^2_\ell$} \\
\cmidrule(lr){3-4} \cmidrule(lr){5-6} \cmidrule(lr){7-8}
$\tau$ & $\rho$ & Lasso & Ridge & Lasso & Ridge & Lasso & Ridge \\
\midrule
0.2 & 0.0 & 9.47 & 6.05 & 49.0 & 60.0 & 0.365 & 0.598 \\
0.1 & 0.0 & 5.85 & 4.11 & 35.0 & 69.0 & 0.130 & 0.427 \\
0.05 & 0.0 & 3.14 & 2.56 & 66.0 & 85.0 & 0.026 & 0.257 \\
0.01 & 0.0 & 0.69 & 0.69 & 92.0 & 92.0 & $-$0.058 & $-$0.033 \\
\midrule
0.1 & 0.5 & 5.91 & 3.59 & 87.0 & 90.0 & 0.742 & 0.844 \\
0.1 & 0.8 & 4.40 & 2.27 & 93.0 & 92.0 & 0.938 & 0.968 \\
\bottomrule
\end{tabular}
\end{table}


\subsubsection*{$\sqrt{n}\cdot\Delta_n$ as an Upper Bound}

An important subtlety emerged in the asymmetric semi-synthetic results (Section~\ref{sec:asymmetric}), where Ridge exhibited higher $\sqrt{n}\cdot\Delta_n$ than Lasso yet achieved equal or better coverage. This apparent contradiction is resolved by recognising that $\sqrt{n}\cdot\Delta_n$ bounds the remainder term via the Cauchy-Schwarz inequality:
\[
|R_n| \;\leq\; \|\hat{\ell} - \ell_0\|_{L_2} \cdot \|\hat{m} - m_0\|_{L_2} \;=\; \Delta_n.
\]
The actual remainder $R_n \approx n^{-1}\sum_i (\hat{\ell}_i - \ell_{0,i})(\hat{m}_i - m_{0,i})$ depends on the \emph{inner product} of the two error vectors, not the product of their norms. When the errors are aligned, as in the symmetric design where $\beta_m = \beta_g$ and both error vectors arise from estimating the same coefficient structure, the bound is tight and $\sqrt{n}\cdot\Delta_n$ tracks coverage reliably. When the errors are misaligned, as in asymmetric designs where $\beta_m$ and $\beta_g$ have different sparsity structures, the bound is loose.

This distinction has a direct implication for practical diagnostics. On real data, neither $\Delta_n$ nor $R_n$ is computable because the true nuisance functions are unknown. The approximation $\hat{\Delta}_n$ based on out-of-fold residuals against observed $Y$ and $D$ inherits both the noise inflation from irreducible $\varepsilon$ and $v$, and the Cauchy-Schwarz looseness. We therefore recommend examining the out-of-fold $R^2$ of each nuisance regression \emph{separately}, comparing Lasso and Ridge. This avoids the product structure entirely and provides two pieces of information that $\sqrt{n}\cdot\Delta_n$ cannot: which nuisance function is driving the failure, and whether the researcher is in the symmetric or asymmetric case.

Table~\ref{tab:cauchy} quantifies this looseness using the FRED-MD semi-synthetic results at $n = 400$, where the approximation $\sqrt{n}\cdot|\text{Bias}|$ serves as a proxy for $\sqrt{n}\cdot|R_n|$.

\begin{table}[H]
\centering
\small
\caption{Cauchy-Schwarz tightness: bound ($\sqrt{n}\cdot\Delta_n$) vs.\ approximate actual remainder ($\sqrt{n}\cdot|\text{Bias}|$) on FRED-MD ($n = 400$). In symmetric dense designs the bound is approximately $5$--$6\times$ the actual remainder. In asymmetric designs, the bound exceeds the actual by two orders of magnitude, explaining why $\sqrt{n}\cdot\Delta_n$ is misleading under asymmetric nuisance structures.}
\label{tab:cauchy}
\begin{tabular}{llcccc}
\toprule
Design & Method & $\sqrt{n}\cdot\Delta_n$ & $\sqrt{n}\cdot|\text{Bias}|$ & Bound/Actual \\
\midrule
Symmetric dense ($\tau = 0.1$) & Lasso & 4.63 & 0.75 & $6\times$ \\
                                & Ridge & 3.27 & 0.60 & $5\times$ \\
\midrule
Asym.\ sparse $\beta_m$ ($\tau_g = 0.1$) & Lasso & 3.17 & 0.02 & $144\times$ \\
                                          & Ridge & 11.03 & 0.05 & $230\times$ \\
\midrule
Asym.\ dense $\beta_m$ ($\tau_m = 0.1$) & Lasso & 4.04 & 0.05 & $88\times$ \\
                                          & Ridge & 5.20 & 0.89 & $6\times$ \\
\midrule
Symmetric sparse ($\tau = 2.0$) & Lasso & 2.57 & 0.06 & $42\times$ \\
                                 & Ridge & 13.89 & 1.39 & $10\times$ \\
\bottomrule
\end{tabular}
\end{table}

This motivates the practical recommendation to compare $R^2$ on each nuisance regression separately rather than relying on the product-based diagnostic.



%% ============================================================
\subsection{What Makes the Breakdown Worse}\label{sec:sensitivity}
%% ============================================================

The 648-cell grid systematically varied five parameters. This section examines how each modulates the coverage gap in the weak dense regime ($\tau \leq 0.2$, $s/p = 1.0$; 107 cells).

\begin{figure}[H]
  \centering
  \includegraphics[width=\textwidth]{fig4.pdf}
  \caption{Parameter sensitivity of the coverage gap (Ridge~$-$~Lasso) in the weak dense regime ($\tau \leq 0.2$, $s/p = 1.0$). (a)~Signal strength: the gap peaks at $\tau = 0.1$ (avg.\ 11.3~pp) and narrows at $\tau = 0.01$ as nuisance functions approach zero. (b)~Dimensionality: higher $p/n$ amplifies the gap from 3.7~pp to 9.7~pp. (c)~Correlation: the gap collapses from 9.3~pp at $\rho = 0.0$ to 1.9~pp at $\rho = 0.8$. (d)~Sample size: the gap widens from 9~pp at $n = 300$ to 34~pp at $n = 1000$ because Lasso's prediction error does not improve with $n$ under dense signals.}
  \label{fig:sensitivity}
\end{figure}

\subsubsection*{Signal Strength ($\tau$)}

The coverage gap was not monotone in $\tau$. Averaging across all cells in the weak dense regime, the gap peaked at $\tau = 0.1$ (mean: 11.3~pp, maximum: 41.0~pp) and declined at both stronger and weaker signals. At $\tau = 0.01$, the mean gap was only 1.1~pp, as both methods converged when nuisance functions approached zero (Figure~\ref{fig:sensitivity}a). The dangerous zone was $\tau \in [0.05, 0.2]$, corresponding to $R^2 \in [5\%, 17\%]$, the range in which most empirical economics papers operate.

\subsubsection*{Dimensionality ($p/n$)}

The average gap grew monotonically with $p/n$: 3.7~pp at $p/n \leq 0.20$, 7.5~pp at $p/n \in (0.20, 0.50]$, and 9.7~pp at $p/n > 0.50$ (Figure~\ref{fig:sensitivity}b). Higher dimensionality increased the number of coefficients for Lasso to incorrectly threshold to zero, amplifying the prediction failure. Both methods deteriorated at high $p/n$, but Ridge deteriorated more slowly: average Ridge coverage was 90.5\% at $p/n \leq 0.20$ versus 76.6\% at $p/n > 0.50$; average Lasso coverage was 86.8\% versus 66.8\%.

\subsubsection*{Predictor Correlation ($\rho$)}

Correlation compressed the gap substantially. At $\rho = 0.0$ and $\rho = 0.5$, the average gaps were 9.3 and 9.8~pp respectively. At $\rho = 0.8$, the gap collapsed to 1.9~pp (Figure~\ref{fig:sensitivity}c). Under high Toeplitz correlation, the eigenvalue decay of the covariance matrix reduced the effective dimensionality, allowing Lasso to find sparse approximations to the dense coefficient vector. This mitigation was not complete: even at $\rho = 0.8$, Ridge's $\sqrt{n}\cdot\Delta_n$ remained lower than Lasso's in 19 of 36 weak dense cells.

\begin{figure}[H]
  \centering
  \includegraphics[width=0.75\textwidth]{fig5.pdf}
  \caption{Interaction between $p/n$ and $\rho$ at $\tau = 0.1$, $s/p = 1.0$, $n = 1000$. The gap at $p/n = 0.50$ widened from 28~pp at $\rho = 0.0$ to 41~pp at $\rho = 0.5$ (non-monotone: Ridge improved faster than Lasso) before collapsing to 1~pp at $\rho = 0.8$.}
  \label{fig:interaction}
\end{figure}

The interaction between $p/n$ and $\rho$ was non-trivial (Figure~\ref{fig:interaction}). At $p/n = 0.50$, the gap at $\rho = 0.5$ (41~pp) exceeded the gap at $\rho = 0.0$ (28~pp). This non-monotonicity arose because intermediate correlation improved Ridge's absolute coverage (from 47\% to 73\%) faster than Lasso's (from 19\% to 32\%). Only at $\rho = 0.8$ did both methods recover sufficiently for the gap to narrow.

\subsubsection*{Sample Size ($n$): More Data Makes It Worse}

The most counterintuitive finding was that increasing $n$ \emph{widened} the coverage gap (Figure~\ref{fig:sensitivity}d). At $p = 200$, $\tau = 0.1$, $s/p = 1.0$, $\rho = 0.0$: the gap was 9~pp at $n = 300$, 27~pp at $n = 500$, and 34~pp at $n = 1000$. This occurred because Lasso's prediction error did not improve with $n$ under dense weak signals, the central theoretical result of \citet{shenxiu2025}. The orthogonality budget $n^{-1/2}$ shrank as $n$ grew, but Lasso's $\Delta_n$ remained roughly constant, so $\sqrt{n}\cdot\Delta_n$ increased: from 3.54 at $n = 300$ to 5.85 at $n = 1000$. Ridge's prediction error did improve with $n$, but not fast enough to offset the shrinking budget entirely. This finding is particularly important for applied researchers who might assume that a large sample protects them from first-stage estimation errors.


\subsubsection*{Two Failure Modes at High Dimensionality}

\begin{figure}[H]
  \centering
  \includegraphics[width=0.8\textwidth]{fig6.pdf}
  \caption{Coverage by $p/n$ ratio at $n = 1000$, $s/p = 1.0$, $\rho = 0.0$. At $p/n = 0.50$ (dotted lines), both methods exhibit near-zero coverage at $\tau \in [0.5, 2.0]$, reflecting a dimensionality failure distinct from the weak-signal failure observed at lower $p/n$. Recovery occurs only as $\tau \to 0$.}
  \label{fig:breakdown}
\end{figure}

Figure~\ref{fig:breakdown} revealed two qualitatively distinct failure modes. At $p/n = 0.10$ and $p/n = 0.20$, the breakdown was signal-driven: coverage dropped as $\tau$ weakened, with Ridge maintaining a consistent advantage. At $p/n = 0.50$ ($n = 1000$, $p = 500$), the pattern was qualitatively different. Both methods exhibited near-zero coverage at strong signals ($\tau = 2.0$: Lasso 7.0\%, Ridge 3.0\%; $\tau = 1.0$: Lasso 0.0\%, Ridge 2.0\%). This was not a weak-signal failure but a dimensionality failure: with $s = p = 500$ active predictors and $n = 1000$ observations, neither method could achieve the convergence rate needed for valid nuisance estimation. Coverage recovered only as $\tau \to 0$ (both 94.0\% at $\tau = 0.01$).


\subsubsection*{Gap Asymmetry: Ridge as a Free Option}

\begin{figure}[H]
  \centering
  \includegraphics[width=\textwidth]{fig7.pdf}
  \caption{Asymmetry of the coverage gap in the weak dense regime. (a)~Ridge wins are large and frequent (75 cells, avg.\ +10.3~pp); Lasso wins are small and rare (15 cells, avg.\ $-$1.6~pp). (b)~When Lasso ``wins,'' both methods already have high coverage, meaning the advantage is negligible. The vertical dashed line in~(b) marks 90\% coverage.}
  \label{fig:asymmetry}
\end{figure}

Figure~\ref{fig:asymmetry} quantifies the asymmetry. Panel~(a) shows that Ridge's wins were both more frequent and larger in magnitude than Lasso's. Panel~(b) shows that the 15 cells where Lasso outperformed Ridge were concentrated at high correlation ($\rho = 0.8$: 10 of 15 cells) and very weak signals ($\tau = 0.01$: 5 of 15 cells), settings where both methods were already at or above 90\% coverage. Of the 15 cells, only one had a gap exceeding 5~pp; the remaining 14 had gaps of 1--2~pp. This confirms that switching to Ridge in the weak dense regime carries no material downside.


\subsubsection*{Ridge Failure in the Sparse Regime}

\begin{figure}[H]
  \centering
  \includegraphics[width=0.7\textwidth]{fig8.pdf}
  \caption{Ridge failure under sparse signals at $n = 1000$, $p = 500$, $s/p = 0.05$, $\rho = 0.0$. Ridge coverage collapsed to 3\% at $\tau = 2.0$ while Lasso maintained 94\%, producing the largest gap in the entire simulation grid (91~pp). Ridge's $\sqrt{n}\cdot\Delta_n$ reached 31.67 versus Lasso's 6.94.}
  \label{fig:ridge_failure}
\end{figure}

While the primary contribution concerns the weak dense regime, the simulation grid also documented the converse failure. At $n = 1000$, $p = 500$, $s/p = 0.05$, $\rho = 0.0$, Ridge coverage was 3.0\% at $\tau = 2.0$ while Lasso maintained 94.0\%, a gap of \textbf{91.0~pp}, the largest in the entire grid (Figure~\ref{fig:ridge_failure}). Eight cells produced exactly 0.0\% Ridge coverage, all at $p \geq n$ with $\rho \leq 0.5$. The mechanism was the well-known inconsistency of Ridge in sparse settings: with only 25 of 500 predictors carrying signal, Ridge applied uniform shrinkage across all 500, pulling the 25 active coefficients toward zero while wasting regularisation capacity on 475 null predictors. This finding prevents the recommendation from being ``always use Ridge'' and motivates the diagnostic approach: compare both methods' $R^2$ on each nuisance regression, and use the method that predicts better.


%% ============================================================
\subsection{When Is the Researcher Protected? Asymmetric Nuisance Structures}\label{sec:asymmetric}
%% ============================================================

The pure simulation used identical coefficient vectors ($\beta_m = \beta_g$) for both nuisance functions. In applied research, the treatment and outcome equations rarely share the same signal structure. The semi-synthetic design tested whether the coverage breakdown persisted when $\beta_m$ and $\beta_g$ had different properties.

\subsubsection*{Sparse Treatment, Dense Weak Outcome}

The first asymmetric design set $\beta_m$ to be sparse and strong ($\tau_m = 2.0$, $s_m/p = 0.05$: 6 of 122 FRED-MD predictors active) and $\beta_g$ to be dense and weak ($\tau_g \in \{0.2, 0.1, 0.05, 0.01\}$, $s_g/p = 1.0$: all 122 active). This models a setting where the treatment assignment depends on a few identifiable variables while the outcome depends diffusely on many indicators.

\begin{table}[H]
\centering
\small
\caption{Asymmetric design on FRED-MD ($n = 400$, $p/n = 0.31$): sparse $\beta_m$ ($\tau_m = 2.0$, $s_m/p = 0.05$) + dense $\beta_g$ ($s_g/p = 1.0$). Lasso outperformed Ridge because $\ell_0$ was dominated by the sparse treatment channel.}
\label{tab:asym_fwd}
\begin{tabular}{lcccccccc}
\toprule
& \multicolumn{2}{c}{Coverage (\%)} & \multicolumn{2}{c}{$\sqrt{n}\cdot\Delta_n$} & \multicolumn{2}{c}{$R^2_\ell$} & \multicolumn{2}{c}{$R^2_m$} \\
\cmidrule(lr){2-3} \cmidrule(lr){4-5} \cmidrule(lr){6-7} \cmidrule(lr){8-9}
$\tau_g$ & Lasso & Ridge & Lasso & Ridge & Lasso & Ridge & Lasso & Ridge \\
\midrule
0.20 & 94.0 & 94.0 & 3.47 & 11.13 & 0.965 & 0.931 & 0.985 & 0.923 \\
0.10 & 95.0 & 94.0 & 3.17 & 11.03 & 0.966 & 0.919 & 0.985 & 0.923 \\
0.05 & 95.0 & 94.0 & 2.92 & 10.98 & 0.967 & 0.909 & 0.985 & 0.923 \\
0.01 & 97.0 & 94.0 & 2.60 & 10.93 & 0.968 & 0.892 & 0.985 & 0.923 \\
\bottomrule
\end{tabular}
\end{table}

Lasso \emph{outperformed} Ridge in every cell (Table~\ref{tab:asym_fwd}). Lasso's $R^2_\ell$ exceeded 0.965 and $R^2_m$ exceeded 0.985 throughout, while Ridge's $\sqrt{n}\cdot\Delta_n$ exceeded 10.9, roughly triple Lasso's values. The explanation lay in the composition of the outcome nuisance function:
\[
\ell_0(X) = \theta_0 \cdot X\beta_m + X\beta_g, \quad \|\beta_m\|^2 = 2.0, \quad \|\beta_g\|^2 \leq 0.2.
\]
The treatment channel $\theta_0 \cdot X\beta_m$ accounted for at least $2.0/(2.0 + 0.2) = 91\%$ of the predictable variation in $\ell_0$. Lasso captured this dominant sparse component with high accuracy, rendering the dense weak $\beta_g$ negligible. Ridge was penalised by applying uniform shrinkage to the sparse treatment equation.

This result also demonstrated the Cauchy-Schwarz looseness discussed in Section~\ref{sec:mechanism}. Ridge had $\sqrt{n}\cdot\Delta_n$ triple that of Lasso yet achieved identical or nearly identical coverage. The product of norms was large, but the actual inner product of the error vectors was small because the two errors were misaligned: Ridge's errors on $\hat{\ell}$ and $\hat{m}$ arose from shrinking a sparse structure (treatment) and a dense structure (outcome) simultaneously, producing error vectors that pointed in different directions.


\subsubsection*{Dense Weak Treatment, Sparse Strong Outcome}

The reverse design ($\beta_m$: $\tau_m \in \{0.1, 0.05\}$, $s_m/p = 1.0$; $\beta_g$: $\tau_g = 2.0$, $s_g/p = 0.05$) forced Lasso to estimate $m_0$ from dense weak signals. Lasso's $R^2_m$ dropped to 0.718 at $n = 400$ with $\tau_m = 0.05$ and to 0.626 at $n = 250$, confirming the prediction failure. Ridge achieved $R^2_m = 0.807$ and 0.770 respectively.

However, coverage held for both methods (Lasso: 94--98\%, Ridge: 92--99\% across all six cells). The mechanism was the same compositional effect operating in the other direction: $\ell_0 = \theta_0 \cdot X\beta_m + X\beta_g$, and with $\|\beta_g\|^2 = 2.0$, the sparse strong $\beta_g$ dominated $\ell_0$. Lasso estimated $\ell_0$ with $R^2_\ell \geq 0.719$, keeping one factor in the product $\Delta_n$ small enough to compensate for the large $m_0$ error.


\subsubsection*{The Central Insight}

The asymmetric results established that the coverage breakdown documented in the symmetric simulation required \emph{both} nuisance functions to involve dense weak signals simultaneously. When either the treatment or the outcome equation contained a dominant sparse component, Lasso exploited the compositional structure $\ell_0 = \theta_0 m_0 + g_0$ and coverage was preserved. This maps directly to applied practice: a researcher whose treatment depends on a few identifiable variables is partially protected, even if the outcome depends diffusely on many weak predictors. The researchers most at risk are those studying settings where both treatment assignment and the outcome are driven by many weak predictors with no dominant sparse component on either side.


%% ============================================================
\subsection{Validation on Real Predictor Data}\label{sec:validation}
%% ============================================================

The pure simulation imposed Toeplitz covariance on the predictors. The semi-synthetic exercises tested whether the mechanism survived under the empirical covariance structures of FRED-MD and Barro-Lee.

\subsubsection*{FRED-MD}

\begin{table}[H]
\centering
\small
\caption{Semi-synthetic FRED-MD: symmetric dense design ($s/p = 1.0$, $\beta_m = \beta_g$). Ridge produced lower $\sqrt{n}\cdot\Delta_n$ in all 12 cells. The Lasso-to-Ridge (L/R) ratio was stable at 1.41--1.60.}
\label{tab:fredmd}
\begin{tabular}{cccccccc}
\toprule
& & \multicolumn{2}{c}{$\sqrt{n}\cdot\Delta_n$} & \multicolumn{2}{c}{Coverage (\%)} & \\
\cmidrule(lr){3-4} \cmidrule(lr){5-6}
$n$ ($p/n$) & $\tau$ & Lasso & Ridge & Lasso & Ridge & L/R ratio \\
\midrule
\multirow{4}{*}{757 (0.16)}
 & 0.20 & 5.75 & 4.08 & 93.0 & 93.0 & 1.41 \\
 & 0.10 & 4.42 & 3.07 & 91.0 & 93.0 & 1.44 \\
 & 0.05 & 3.22 & 2.26 & 93.0 & 95.0 & 1.43 \\
 & 0.01 & 1.58 & 1.05 & 91.0 & 95.0 & 1.50 \\
\midrule
\multirow{4}{*}{400 (0.31)}
 & 0.20 & 6.46 & 4.54 & \textbf{81.0} & \textbf{90.0} & 1.42 \\
 & 0.10 & 4.63 & 3.27 & 85.0 & 90.0 & 1.42 \\
 & 0.05 & 3.43 & 2.34 & 89.0 & 94.0 & 1.47 \\
 & 0.01 & 1.71 & 1.07 & 93.0 & 93.0 & 1.60 \\
\midrule
\multirow{4}{*}{250 (0.49)}
 & 0.20 & 7.72 & 5.34 & 92.0 & 94.0 & 1.45 \\
 & 0.10 & 5.47 & 3.67 & 97.0 & 94.0 & 1.49 \\
 & 0.05 & 3.97 & 2.55 & 98.0 & 95.0 & 1.56 \\
 & 0.01 & 1.76 & 1.13 & 96.0 & 95.0 & 1.56 \\
\bottomrule
\end{tabular}
\end{table}

The mechanism replicated in full (Table~\ref{tab:fredmd}). Ridge produced a strictly lower $\sqrt{n}\cdot\Delta_n$ than Lasso in all 12 symmetric dense cells, with the L/R ratio ranging from 1.41 to 1.60. The largest coverage gap was 9~pp at $n = 400$, $\tau = 0.2$ (Lasso: 81.0\%, Ridge: 90.0\%). Coverage gaps were smaller than the pure simulation (where they reached 34~pp at comparable $p/n$) because $p = 122$ is smaller than $p = 200$--$500$, limiting the number of coefficients Lasso incorrectly thresholded.

\begin{figure}[H]
  \centering
  \includegraphics[width=\textwidth]{fig9.pdf}
  \caption{$\sqrt{n}\cdot\Delta_n$ on FRED-MD across three sample sizes ($s/p = 1.0$, symmetric dense). Ridge produced strictly lower values in all 12 cells. The gap widened with $p/n$, consistent with the dimensionality scaling documented in the pure simulation.}
  \label{fig:fredmd_sqrtNDn}
\end{figure}

\begin{figure}[H]
  \centering
  \includegraphics[width=\textwidth]{fig10.pdf}
  \caption{Coverage on FRED-MD across three sample sizes. The largest gap (9~pp) occurred at $n = 400$, $\tau = 0.2$. At $n = 250$, Lasso exhibited inflated coverage (97--98\%) due to wider confidence intervals from estimation imprecision, not superior nuisance estimation.}
  \label{fig:fredmd_coverage}
\end{figure}

At $n = 250$ ($p/n = 0.49$), Lasso exhibited unexpectedly high coverage of 97--98\% (Figure~\ref{fig:fredmd_coverage}). This was not evidence of superior performance. Lasso's RMSE at $n = 250$ ranged from 0.064 to 0.112, substantially larger than at $n = 400$ (0.057--0.081), producing wider confidence intervals that covered the true value more frequently despite the point estimate being less precise. Ridge's coverage at $n = 250$ remained at 94--95\%, indicating better-calibrated intervals.


\subsubsection*{Barro-Lee}

The Barro-Lee dataset provided the highest dimensionality ratio ($p/n = 0.678$) and a qualitatively different correlation structure (mean $|\rho_{jk}| = 0.452$). The high correlation would typically compress the coverage gap; the high $p/n$ would amplify it.

\begin{table}[H]
\centering
\small
\caption{Semi-synthetic Barro-Lee: symmetric dense design ($n = 90$, $p = 61$, $p/n = 0.678$, $B = 50$). Ridge maintained lower $\sqrt{n}\cdot\Delta_n$ in all five cells despite the high predictor correlation. At $\tau = 0.01$, Lasso's $R^2_m = -0.035$ indicates performance worse than a constant predictor.}
\label{tab:barrolee}
\begin{tabular}{ccccccccc}
\toprule
& \multicolumn{2}{c}{Coverage (\%)} & \multicolumn{2}{c}{$\sqrt{n}\cdot\Delta_n$} & \multicolumn{2}{c}{$R^2_\ell$} & \multicolumn{2}{c}{$R^2_m$} \\
\cmidrule(lr){2-3} \cmidrule(lr){4-5} \cmidrule(lr){6-7} \cmidrule(lr){8-9}
$\tau$ & Lasso & Ridge & Lasso & Ridge & Lasso & Ridge & Lasso & Ridge \\
\midrule
0.50 & 90.0 & 92.0 & 5.69 & 3.53 & 0.949 & 0.969 & 0.922 & 0.951 \\
0.20 & 90.0 & 94.0 & 3.84 & 2.55 & 0.915 & 0.944 & 0.865 & 0.909 \\
0.10 & \textbf{86.0} & \textbf{94.0} & 2.92 & 2.02 & 0.872 & 0.911 & 0.791 & 0.855 \\
0.05 & 86.0 & 94.0 & 2.33 & 1.62 & 0.799 & 0.859 & 0.658 & 0.763 \\
0.01 & 90.0 & 92.0 & 1.45 & 1.07 & 0.353 & 0.538 & $-$0.035 & 0.178 \\
\bottomrule
\end{tabular}
\end{table}

\begin{figure}[H]
  \centering
  \includegraphics[width=\textwidth]{fig11.pdf}
  \caption{Semi-synthetic Barro-Lee ($n = 90$, $p = 61$, $p/n = 0.678$). (a)~$\sqrt{n}\cdot\Delta_n$: Ridge maintained lower values throughout. (b)~Coverage: the gap reached 8~pp at $\tau = 0.1$ and $\tau = 0.05$ (Lasso: 86\%, Ridge: 94\%) despite the high predictor correlation (mean $|\rho_{jk}| = 0.452$).}
  \label{fig:barrolee}
\end{figure}

The coverage gap reached 8~pp at $\tau = 0.1$ and $\tau = 0.05$ (Lasso: 86\%, Ridge: 94\%), despite the high predictor correlation (Table~\ref{tab:barrolee}, Figure~\ref{fig:barrolee}). Ridge produced lower $\sqrt{n}\cdot\Delta_n$ in all five symmetric dense cells, with L/R ratios of 1.35--1.61. At $\tau = 0.01$, Lasso's $R^2_m$ dropped to $-0.035$, confirming that Lasso performed worse than a constant predictor on this dataset under very weak dense signals.

The Barro-Lee sparse baseline confirmed Ridge failure at this dimensionality: coverage fell to 78.0\% for Ridge (Lasso: 90.0\%), with Ridge bias of 0.111 and $\sqrt{n}\cdot\Delta_n = 8.10$ (Lasso: 4.58).

\subsubsection*{Cross-Dataset Consistency}

\begin{figure}[H]
  \centering
  \includegraphics[width=0.82\textwidth]{fig12.pdf}
  \caption{$\sqrt{n}\cdot\Delta_n$ at $\tau = 0.1$, symmetric dense, across all datasets. Ridge produced lower values in every configuration. The Lasso-to-Ridge ratio (annotated above each pair) was stable at 1.42--1.49, confirming that the mechanism is robust to the covariance structure.}
  \label{fig:cross_dataset}
\end{figure}

Figure~\ref{fig:cross_dataset} compares the nuisance error product across all datasets at $\tau = 0.1$. The Lasso-to-Ridge ratio was remarkably stable: 1.42 in the pure simulation ($p/n = 0.20$, $\rho = 0$), 1.44 on FRED-MD ($p/n = 0.16$), 1.42 on FRED-MD ($p/n = 0.31$), 1.49 on FRED-MD ($p/n = 0.49$), and 1.45 on Barro-Lee ($p/n = 0.68$). This consistency across Toeplitz and empirical covariance structures, across macro and growth contexts, and across $p/n$ from 0.16 to 0.68, confirmed that the mechanism was not an artefact of any particular simulation design choice.


\subsection{The Diagnostic in Practice}\label{sec:empirical_results}

\textit{This section will present the empirical diagnostic results from the Barro-Lee and Goyal-Welch datasets, applying the four-step diagnostic procedure described in Section~\ref{sec:empirical_design}. For each dataset and treatment variable, the diagnostic reports: (1)~signal strength positioning via Ridge $R^2$, (2)~separate $R^2$ comparison between Lasso and Ridge on each nuisance regression, (3)~approximate $\sqrt{n}\cdot\hat{\Delta}_n$, and (4)~estimate stability across cross-fitting splits. Preliminary results are available and will be integrated upon completion of the robustness checks.}



%% ============================================================
\section{Practical Implications}\label{sec:implications}
%% ============================================================

The results have three implications for applied researchers implementing double machine learning in high-dimensional settings. Taken together, they suggest that nuisance learner choice should be treated as an empirical modelling decision rather than as a routine default.

\medskip
\noindent\textbf{Recommendation 1. Lasso should not be treated as the default nuisance learner.}
The approximate sparsity condition that underpins Lasso is difficult to reconcile with an accumulating body of evidence indicating that economic signals are often diffuse and distributed across many predictors rather than concentrated in a small subset \citep{giannone2021, gu2020}. In our simulations, at signal strengths comparable to those frequently encountered in applied economics ($R^2 \approx 9\%$, corresponding to $\tau = 0.1$), coverage for Lasso-based DML fell to 35\%. This is not a modest deterioration in finite-sample performance. Rather, it implies that the nominal 95\% confidence interval excludes the true parameter in 65\% of repeated samples. In such environments, conventional use of DML-Lasso can therefore produce a false appearance of precision and induce researchers to report statistically significant effects that are not warranted by the underlying data-generating process. These findings suggest that the widespread practice of treating Lasso as the default nuisance learner in applied DML should be reconsidered.

\medskip
\noindent\textbf{Recommendation 2. Compare out-of-fold $R^2$ separately for the two nuisance regressions.}
Before selecting a nuisance learner, researchers should estimate both nuisance functions using Lasso and Ridge and compare the resulting out-of-fold $R^2$ values separately for the outcome and treatment regressions. The relevant comparison is thus at the level of each nuisance equation, rather than through a single product-based summary statistic.

If Ridge achieves substantially higher out-of-fold $R^2$ than Lasso in both nuisance regressions, the empirical environment is likely characterised by dense weak signals on both sides of the partially linear model. This is the symmetric regime in which the coverage deterioration documented in the preceding sections is most severe. In such cases, Ridge should generally be preferred.

If the $R^2$ advantage of Ridge arises in only one nuisance regression, the setting is asymmetric. The semi-synthetic results indicate that when either the treatment or the outcome equation contains a dominant sparse component, Lasso can still exploit the compositional structure of the outcome nuisance function, $\ell_0 = \theta_0 m_0 + g_0$, thereby partially preserving inferential validity. In this case Ridge remains a relevant alternative, but the inferential risk associated with Lasso is less acute than in the fully symmetric dense regime.

If Lasso and Ridge produce similar out-of-fold $R^2$ values in both nuisance regressions, two interpretations are plausible. Either the signal structure is sufficiently sparse for Lasso to perform well, or the signals are strong enough that both learners satisfy the rate requirements needed for reliable inference. In either case, the choice between the two methods is unlikely to materially affect the resulting confidence intervals.

This diagnostic has two advantages. First, it is feasible in practice because it does not require knowledge of the true nuisance functions. Second, it avoids the interpretive difficulties associated with the product-based quantity $\sqrt{n}\cdot\Delta_n$. As shown in Section~\ref{sec:mechanism}, $\sqrt{n}\cdot\Delta_n$ is a Cauchy--Schwarz upper bound on the scaled remainder. That bound is informative in symmetric settings, but it can be highly non-tight in asymmetric designs. In our results (Table~\ref{tab:cauchy}), the ratio of the bound to the realised remainder was approximately 6 in symmetric designs but as high as 230 in asymmetric ones. By contrast, separate out-of-fold $R^2$ comparisons are not affected by this looseness and additionally identify which nuisance regression is driving the inferential vulnerability.

Large differences between the causal estimates $\hat{\theta}$ obtained from DML-Lasso and DML-Ridge are themselves substantively informative. Such divergence indicates that the estimated treatment effect is sensitive to the choice of nuisance learner, which in turn suggests that the underlying rate condition may be close to failure. In those circumstances, reporting both estimates and discussing their sensitivity is more informative than selecting a single specification without comment.

\medskip
\noindent\textbf{Recommendation 3. Pay close attention to dimensionality.}
The inferential consequences of nuisance learner choice become more pronounced as the dimensionality ratio $p/n$ increases. In our simulations, when $p/n \leq 0.20$, the average coverage gap between methods was 3.7 percentage points, whereas for $p/n > 0.50$ it increased to 9.7 percentage points. More importantly, at $p/n = 0.50$ under dense signals and strong predictive strength, both Lasso and Ridge exhibited near-zero coverage, indicating that neither estimator may be adequate in sufficiently high-dimensional environments. For applied researchers working with datasets in which $p/n$ exceeds 0.50 and the signal structure is plausibly dense, the relevant question is not simply whether to use Lasso or Ridge, but whether linear penalised regression is an adequate first-stage strategy at all. In such settings, dimension-reduction methods, including principal components or factor-augmented approaches, or alternative nonlinear learners may be more appropriate.

A related implication concerns sample size. One might expect that larger samples automatically mitigate the effect of first-stage estimation error on DML inference. Our results show that this intuition can be misleading under dense weak signals. In that regime, the orthogonality budget shrinks at the rate $n^{-1/2}$, while Lasso's prediction error improves too slowly to offset this contraction. As a result, $\sqrt{n}\cdot\Delta_n$ can increase with $n$ rather than decrease. A researcher working with $n = 1000$ may therefore be more exposed to inferential failure than one working with $n = 300$, despite the larger sample. Large samples do not, by themselves, guarantee protection from poor nuisance estimation when the signal environment is unfavourable.

\medskip
These recommendations should be interpreted within the scope of the present analysis. The paper examines the choice between Lasso and Ridge as linear nuisance learners under alternative signal structures; it does not provide a general ranking of all possible first-stage methods. In particular, elastic net represents a natural intermediate case that is not studied here. Likewise, nonlinear methods such as random forests, gradient boosting, and neural networks may respond differently to dense weak signals and remain important topics for future research. The implications developed here are therefore most directly applicable to the linear, moderate- to high-dimensional settings ($p/n \in [0.10, 0.70]$) that characterise much of the applied DML work in economics.


%% ============================================================
\section{Conclusion}\label{sec:conclusion}
%% ============================================================

This paper investigated whether the prediction failure of Lasso under dense weak signals, documented by \citet{shenxiu2025}, compromises causal inference in Double Machine Learning. We framed the question through the concept of orthogonality as a finite insurance mechanism: the DML estimator absorbs nuisance estimation errors up to a budget of size $n^{-1/2}$, and inference fails when the nuisance error product $\Delta_n$ exceeds this budget.

Three main findings emerged from a 648-cell Monte Carlo simulation, semi-synthetic validation on FRED-MD and Barro-Lee predictor data, and asymmetric nuisance designs.

First, Lasso-based DML produces invalid confidence intervals under dense weak signals. Coverage dropped to 35\% at signal strengths corresponding to the 25th percentile of published economics papers, while Ridge-based DML maintained 69\%. The mechanism traces cleanly from prediction failure through an elevated nuisance error product to coverage collapse. Counterintuitively, increasing sample size widened the gap rather than narrowing it, because Lasso's prediction error does not improve with $n$ under dense signals while the orthogonality budget shrinks. This is a direct downstream consequence of the theoretical result of \citet{shenxiu2025}, applied for the first time to the causal inference setting.

Second, the coverage breakdown requires both nuisance functions to involve dense weak signals simultaneously. When either the treatment or the outcome equation contains a dominant sparse component, Lasso exploits the compositional structure of the outcome nuisance function $\ell_0 = \theta_0 m_0 + g_0$ and coverage is preserved. This finding narrows the scope of the vulnerability: the researchers most at risk are those studying settings where both treatment assignment and the outcome depend on many weak predictors, with no dominant sparse component on either side.

Third, neither method dominates universally. Ridge failed catastrophically under sparse strong signals, with coverage dropping to 3\% (versus Lasso's 94\%) in the sparse high-dimensional regime. The practical contribution is therefore not to recommend Ridge unconditionally, but to provide a diagnostic: comparing out-of-fold $R^2$ from both methods on each nuisance regression separately identifies the researcher's exposure and guides the choice. This diagnostic avoids the ambiguities of the product-based $\sqrt{n}\cdot\Delta_n$ measure, which we showed is a Cauchy-Schwarz upper bound that can be loose by two orders of magnitude under asymmetric nuisance structures (Table~\ref{tab:cauchy}).

The semi-synthetic validation confirmed the mechanism on real predictor data. The Lasso-to-Ridge $\sqrt{n}\cdot\Delta_n$ ratio was stable at 1.41 to 1.61 across FRED-MD and Barro-Lee, matching the pure simulation despite markedly different covariance structures. On Barro-Lee ($p/n = 0.678$, mean $|\rho_{jk}| = 0.452$), the coverage gap reached 8 percentage points even under the high predictor correlation that the pure simulation showed typically compresses the gap. These results confirm that the mechanism is not an artefact of the Toeplitz covariance assumed in the pure simulation.


\subsection*{Limitations}

Several limitations should be acknowledged. The data-generating process in both the pure simulation and the semi-synthetic exercises assumed linear nuisance functions ($m_0(X) = X\beta_m$, $g_0(X) = X\beta_g$). Real nuisance functions may be nonlinear, and the relative performance of Lasso and Ridge could differ in the presence of interactions, threshold effects, or other nonlinearities. The noise terms were drawn from Gaussian distributions with unit variance; heavy-tailed or heteroskedastic errors could alter the coverage properties of both methods.

The coefficient construction assigned equal magnitude $\sqrt{\tau/s}$ to all active predictors. In practice, coefficients are heterogeneous, and Lasso may perform better when a subset of predictors has substantially larger coefficients than the rest, even if the overall structure is dense. This represents a setting intermediate between the sparse and dense extremes studied here.

The replication count of $B = 100$ ($B = 50$ for Barro-Lee) was adequate for detecting the large coverage differences documented in the paper but produces coverage estimates with standard errors of approximately 2.2~pp (3.1~pp for Barro-Lee). Differences below 5~pp should therefore be interpreted with caution.

The paper evaluated only two nuisance learners: Lasso and Ridge. Elastic net, which combines $\ell_1$ and $\ell_2$ penalties, represents a natural intermediate that could potentially capture the benefits of both methods. Random forests, gradient boosting, and neural networks may respond differently to dense weak signals. Extending the comparison to these methods is a natural direction for future work.

Finally, the Toeplitz covariance structure used in the pure simulation, while a standard choice, imposes a symmetric and smoothly decaying dependence pattern. The semi-synthetic exercises partially address this limitation by using real predictor matrices, but the semi-synthetic designs still impose linear nuisance functions on these matrices. A fully nonparametric simulation study, while computationally demanding, would provide stronger evidence on the generalisability of the findings.


\subsection*{Directions for Future Research}

Three extensions appear most promising. First, a theoretical analysis formalising the conditions under which Lasso's nuisance error product diverges while Ridge's converges would provide the formal underpinning for the simulation evidence presented here. Second, extending the comparison to elastic net and nonlinear learners would test whether the dense weak-signal vulnerability is specific to Lasso or shared by other sparsity-inducing methods. Third, applying the diagnostic to real causal applications where the treatment effect is debated, such as the effect of monetary policy on inflation or the returns to education, would demonstrate the practical consequences of the nuisance learner choice in settings where the true answer is unknown and the stakes are high.

The central message of this paper is that the choice of nuisance learner in DML is not a technical detail. It is a substantive decision that determines whether the causal estimates a researcher reports can be trusted. Orthogonality provides protection, but that protection has a limit. Under the dense weak signals that characterise much of economic data, Lasso exhausts that limit while Ridge does not.


\newpage
%% ---- References --------------------------------------------
\bibliographystyle{apalike}
\bibliography{MyReferencesFile}


%% ============================================================
%% APPENDICES
%% ============================================================
\appendix

\section{DML Algorithm and Implementation Details}\label{app:algorithm}

Algorithm~\ref{alg:dml} describes the full DML procedure as implemented in this paper.

\begin{algorithm}[H]
\caption{DML Estimation with Cross-Fitting}\label{alg:dml}
\begin{algorithmic}[1]
\For{$b = 1, \ldots, B$} \Comment{Monte Carlo replications}
    \State Draw $X^{(b)}, v^{(b)}, \varepsilon^{(b)}$; compute $D^{(b)}, Y^{(b)}$
    \For{$s = 1, \ldots, S$} \Comment{Cross-fitting repetitions}
        \State Partition $\{1, \ldots, n\}$ into $K$ folds $\mathcal{I}_1^{(s)}, \ldots, \mathcal{I}_K^{(s)}$
        \For{$k = 1, \ldots, K$}
            \State $\mathcal{I}^c_k \gets \{1, \ldots, n\} \setminus \mathcal{I}_k^{(s)}$
            \State Standardise $X_{\mathcal{I}^c_k}$ to mean 0, variance 1; apply same transform to $X_{\mathcal{I}_k}$
            \State Fit $\hat{\ell}^{(k)}$ on $(\tilde{X}_{\mathcal{I}^c_k}, Y_{\mathcal{I}^c_k})$ using method $\in \{\text{LassoCV}, \text{RidgeCV}\}$
            \State Fit $\hat{m}^{(k)}$ on $(\tilde{X}_{\mathcal{I}^c_k}, D_{\mathcal{I}^c_k})$ using same method
            \State Compute $\hat{W}_i = Y_i - \hat{\ell}^{(k)}(X_i)$, $\hat{V}_i = D_i - \hat{m}^{(k)}(X_i)$ for $i \in \mathcal{I}_k$
        \EndFor
        \State $\hat{\theta}^{(s)} \gets \sum_i \hat{V}_i \hat{W}_i / \sum_i \hat{V}_i^2$
        \State $\hat{\sigma}^{(s)2} \gets$ sandwich formula~\eqref{eq:variance}
    \EndFor
    \State $\tilde{\theta}_b \gets S^{-1}\sum_s \hat{\theta}^{(s)}$; $\tilde{\sigma}^2_b \gets S^{-1}\sum_s[\hat{\sigma}^{(s)2} + (\hat{\theta}^{(s)} - \tilde{\theta}_b)^2]$
    \State $\text{CI}_b \gets \tilde{\theta}_b \pm 1.96 \cdot \tilde{\sigma}_b$
\EndFor
\State Report: Bias, RMSE, Coverage, $\sqrt{n}\cdot\Delta_n$, $R^2_\ell$, $R^2_m$
\end{algorithmic}
\end{algorithm}

\noindent\textbf{LassoCV:} $\lambda$ selected over 15 log-spaced values from $10^{-4}$ to $10^1$ via $K_{\text{inner}} = 3$-fold CV, minimising MSE. Implementation: \texttt{sklearn.linear\_model.LassoCV}.

\noindent\textbf{RidgeCV:} $\lambda$ selected over 60 log-spaced values from $10^{-4}$ to $10^6$ via $K_{\text{inner}} = 3$-fold CV, scoring by negative MSE. The wider grid reflects the broader range of effective regularisation across the signal spectrum. Implementation: \texttt{sklearn.linear\_model.RidgeCV}.

Cross-validated $\lambda$ selection was chosen over theoretical penalties (e.g., BRT; \citealt{belloni2014}) because cross-validation is the dominant approach in applied DML practice, and the goal of this study was to evaluate the methods as they are actually used, not under idealised theoretical conditions.


\section{FRED-MD Preprocessing}\label{app:fredmd}

The FRED-MD dataset \citep{mccracken2016} was preprocessed as follows. The raw dataset contained 126 series spanning January 1959 to January 2026. Each series was transformed to stationarity using the McCracken-Ng transformation codes: no change (code 1, 11 series), first differences (code 2, 19 series), log levels (code 4, 10 series), log first differences (code 5, 52 series), log second differences (code 6, 33 series), and first differences of percentage changes (code 7, 1 series). Columns with more than 10\% missing values were dropped (4 series removed), and rows with any remaining missing observations were removed, yielding a final predictor matrix of $n = 757$ observations and $p = 122$ predictors ($p/n = 0.161$). All columns were standardised to mean zero and unit variance. Subsamples ($n = 400$, $n = 250$) were drawn without replacement using a fixed random seed and re-standardised.


\section{Full Semi-Synthetic Results}\label{app:semi_results}

Tables~\ref{tab:app_fredmd_asym} and~\ref{tab:app_barrolee_full} present the complete semi-synthetic results not shown in the main text.

\begin{table}[H]
\centering
\small
\caption{FRED-MD asymmetric designs and sparse baseline (all sample sizes). Asym.~A: sparse $\beta_m$ ($\tau_m = 2.0$, $s_m/p = 0.05$) + dense $\beta_g$. Asym.~B: dense $\beta_m$ + sparse $\beta_g$ ($\tau_g = 2.0$, $s_g/p = 0.05$). Sparse: $\beta_m = \beta_g$, $s/p = 0.05$, $\tau = 2.0$.}
\label{tab:app_fredmd_asym}
\begin{tabular}{llccccccc}
\toprule
& & \multicolumn{2}{c}{Coverage (\%)} & \multicolumn{2}{c}{$\sqrt{n}\cdot\Delta_n$} & \multicolumn{2}{c}{$R^2_\ell$} \\
\cmidrule(lr){3-4} \cmidrule(lr){5-6} \cmidrule(lr){7-8}
$n$ & Design ($\tau_g$ or $\tau_m$) & Lasso & Ridge & Lasso & Ridge & Lasso & Ridge \\
\midrule
\multicolumn{8}{l}{\textit{Asymmetric A: sparse $\beta_m$ + dense $\beta_g$}} \\
757 & $\tau_g = 0.10$ & 95.0 & 93.0 & 2.77 & 9.83 & 0.975 & 0.943 \\
400 & $\tau_g = 0.10$ & 95.0 & 94.0 & 3.17 & 11.03 & 0.966 & 0.919 \\
250 & $\tau_g = 0.10$ & 95.0 & 93.0 & 3.77 & 12.32 & 0.951 & 0.887 \\
\midrule
\multicolumn{8}{l}{\textit{Asymmetric B: dense $\beta_m$ + sparse $\beta_g$}} \\
757 & $\tau_m = 0.10$ & 95.0 & 96.0 & 3.60 & 4.47 & 0.717 & 0.798 \\
400 & $\tau_m = 0.10$ & 94.0 & 99.0 & 4.04 & 5.20 & 0.719 & 0.793 \\
250 & $\tau_m = 0.10$ & 98.0 & 95.0 & 4.76 & 6.24 & 0.715 & 0.776 \\
\midrule
\multicolumn{8}{l}{\textit{Symmetric sparse baseline}} \\
757 & $\tau = 2.0$ & 95.0 & 93.0 & 1.76 & 10.29 & 0.986 & 0.934 \\
400 & $\tau = 2.0$ & 94.0 & 93.0 & 2.57 & 13.89 & 0.976 & 0.900 \\
250 & $\tau = 2.0$ & 90.0 & 80.0 & 4.58 & 18.24 & 0.955 & 0.838 \\
\bottomrule
\end{tabular}
\end{table}

\begin{table}[H]
\centering
\small
\caption{Barro-Lee full semi-synthetic results ($n = 90$, $p = 61$, $B = 50$). Dense designs ($s/p = 1.0$), asymmetric designs, and sparse baseline ($s/p = 0.05$, $\tau = 2.0$).}
\label{tab:app_barrolee_full}
\begin{tabular}{llcccccc}
\toprule
& & \multicolumn{2}{c}{Coverage (\%)} & \multicolumn{2}{c}{$\sqrt{n}\cdot\Delta_n$} & \multicolumn{2}{c}{Bias} \\
\cmidrule(lr){3-4} \cmidrule(lr){5-6} \cmidrule(lr){7-8}
Design & $\tau$ & Lasso & Ridge & Lasso & Ridge & Lasso & Ridge \\
\midrule
\multicolumn{8}{l}{\textit{Symmetric dense ($s/p = 1.0$)}} \\
 & 0.50 & 90.0 & 92.0 & 5.69 & 3.53 & 0.054 & 0.031 \\
 & 0.20 & 90.0 & 94.0 & 3.84 & 2.55 & 0.045 & 0.030 \\
 & 0.10 & 86.0 & 94.0 & 2.92 & 2.02 & 0.037 & 0.026 \\
 & 0.05 & 86.0 & 94.0 & 2.33 & 1.62 & 0.027 & 0.019 \\
 & 0.01 & 90.0 & 92.0 & 1.45 & 1.07 & 0.012 & 0.010 \\
\midrule
\multicolumn{8}{l}{\textit{Asymmetric: sparse $\beta_m$ + dense $\beta_g$ ($\tau_m = 2.0$, $\tau_g = 0.1$)}} \\
 & --- & 92.0 & 86.0 & 2.50 & 6.53 & 0.021 & $-$0.004 \\
\midrule
\multicolumn{8}{l}{\textit{Asymmetric: dense $\beta_m$ + sparse $\beta_g$ ($\tau_m = 0.1$, $\tau_g = 2.0$)}} \\
 & --- & 96.0 & 92.0 & 2.80 & 3.75 & $-$0.009 & 0.045 \\
\midrule
\multicolumn{8}{l}{\textit{Symmetric sparse ($s/p = 0.05$, $\tau = 2.0$)}} \\
 & 2.0 & 90.0 & 78.0 & 4.58 & 8.10 & 0.032 & 0.111 \\
\bottomrule
\end{tabular}
\end{table}


\section{Code Availability}\label{app:code}

All simulation code is available in the project repository. The main scripts are:

\begin{itemize}[leftmargin=*]
    \item \texttt{simulation.py}: 648-cell pure simulation. Parallelised via \texttt{joblib}. Checkpoint system saves each cell to CSV.
    \item \texttt{semi\_synthetic.py}: FRED-MD semi-synthetic (33 cells, $B = 100$).
    \item \texttt{semi\_synthetic\_barrolee.py}: Barro-Lee semi-synthetic (11 cells, $B = 50$).
    \item \texttt{empirical.py}: Empirical diagnostics on Barro-Lee and Goyal-Welch.
\end{itemize}

\noindent Total runtime: approximately 3,804 minutes on an 8-core Apple M-series MacBook Pro (pure simulation) plus 10 minutes (Barro-Lee semi-synthetic) and 60 minutes (FRED-MD semi-synthetic).


\end{document}
